\documentclass[11pt]{article}
\usepackage{amsmath}
\usepackage{tcolorbox}
% Change "review" to "final" to generate the final (sometimes called camera-ready) version.
% Change to "preprint" to generate a non-anonymous version with page numbers.
\usepackage[review]{acl}
 \usepackage{multirow}

% Standard package includes
\usepackage{times}
\usepackage{latexsym}

% For proper rendering and hyphenation of words containing Latin characters (including in bib files)
\usepackage[T1]{fontenc}
% For Vietnamese characters
% \usepackage[T5]{fontenc}
% See https://www.latex-project.org/help/documentation/encguide.pdf for other character sets

% This assumes your files are encoded as UTF8
\usepackage[utf8]{inputenc}
\newcommand{\todoplace}[1]{\textbf{\textcolor{red}{#1}}}
% This is not strictly necessary, and may be commented out,
% but it will improve the layout of the manuscript,
% and will typically save some space.
\usepackage{microtype}
\usepackage{booktabs}

% This is also not strictly necessary, and may be commented out.
% However, it will improve the aesthetics of text in
% the typewriter font.
\usepackage{inconsolata}

%Including images in your LaTeX document requires adding
%additional package(s)
\usepackage{graphicx}

% If the title and author information does not fit in the area allocated, uncomment the following
%
%\setlength\titlebox{<dim>}
%
% and set <dim> to something 5cm or larger.

\title{HiReKG: High-Recall Knowledge Graph Extraction from Plain Text with Language Models}

% Author information can be set in various styles:
% For several authors from the same institution:
% \author{Author 1 \and ... \and Author n \\
%         Address line \\ ... \\ Address line}
% if the names do not fit well on one line use
%         Author 1 \\ {\bf Author 2} \\ ... \\ {\bf Author n} \\
% For authors from different institutions:
% \author{Author 1 \\ Address line \\  ... \\ Address line
%         \And  ... \And
%         Author n \\ Address line \\ ... \\ Address line}
% To start a separate ``row'' of authors use \AND, as in
% \author{Author 1 \\ Address line \\  ... \\ Address line
%         \AND
%         Author 2 \\ Address line \\ ... \\ Address line \And
%         Author 3 \\ Address line \\ ... \\ Address line}

% \author{First Author \\
%   Affiliation / Address line 1 \\
%   Affiliation / Address line 2 \\
%   Affiliation / Address line 3 \\
%   \texttt{email@domain} \\\And
%   Second Author \\
%   Affiliation / Address line 1 \\
%   Affiliation / Address line 2 \\
%   Affiliation / Address line 3 \\
%   \texttt{email@domain} \\}

%\author{
%  \textbf{First Author\textsuperscript{1}},
%  \textbf{Second Author\textsuperscript{1,2}},
%  \textbf{Third T. Author\textsuperscript{1}},
%  \textbf{Fourth Author\textsuperscript{1}},
%\\
%  \textbf{Fifth Author\textsuperscript{1,2}},
%  \textbf{Sixth Author\textsuperscript{1}},
%  \textbf{Seventh Author\textsuperscript{1}},
%  \textbf{Eighth Author \textsuperscript{1,2,3,4}},
%\\
%  \textbf{Ninth Author\textsuperscript{1}},
%  \textbf{Tenth Author\textsuperscript{1}},
%  \textbf{Eleventh E. Author\textsuperscript{1,2,3,4,5}},
%  \textbf{Twelfth Author\textsuperscript{1}},
%\\
%  \textbf{Thirteenth Author\textsuperscript{3}},
%  \textbf{Fourteenth F. Author\textsuperscript{2,4}},
%  \textbf{Fifteenth Author\textsuperscript{1}},
%  \textbf{Sixteenth Author\textsuperscript{1}},
%\\
%  \textbf{Seventeenth S. Author\textsuperscript{4,5}},
%  \textbf{Eighteenth Author\textsuperscript{3,4}},
%  \textbf{Nineteenth N. Author\textsuperscript{2,5}},
%  \textbf{Twentieth Author\textsuperscript{1}}
%\\
%\\
%  \textsuperscript{1}Affiliation 1,
%  \textsuperscript{2}Affiliation 2,
%  \textsuperscript{3}Affiliation 3,
%  \textsuperscript{4}Affiliation 4,
%  \textsuperscript{5}Affiliation 5
%\\
%  \small{
%    \textbf{Correspondence:} \href{mailto:email@domain}{email@domain}
%  }
%}

\begin{document}
\maketitle 
\begin{abstract}
    Knowledge graphs (KGs) represent information as an interlinked network of entities, capturing the relational context that remains implicit in unstructured text. Most existing LLM-based text to KG converters follow an \emph{entity-first} pipeline: entities are extracted upfront, and relations are derived only between them. This design favors precision but systematically misses facts about a single entity, capping recall. We propose \textbf{HiReKG}, a high-recall extraction framework that reduces reliance on upfront entity extraction and introduces an object resolution module to keep the resulting triples cohesive. We evaluate HiReKG across three diverse datasets spanning general, scientific, and technical domains, using three LLM backbones. HiReKG improves recall over the strongest baseline by 18 percentage points on average (up to 29), with a precision trade-off of 2--9 points (5.7 on average), and recovers the single-entity facts that entity-first pipelines structurally miss.
\end{abstract}

\section{Introduction}
\label{sec:introduction}

Knowledge graphs (KGs) provide a structured representation of entities and the relations between them, enabling explicit multi-hop traversal that unstructured text alone does not support. Integrating KGs into retrieval-augmented generation and recommendation systems has been shown to substantially improve their ability to consolidate scattered context from source data \citep{llm-kg-emnlp-survey, recommendation_kg, text_summarization_kg}. These benefits have motivated a long line of work on automatic text-to-KG extraction, from classical OpenIE \citep{openie} to recent LLM-based approaches \citep{kg-llm-survey-arxiv}, all aimed at converting unstructured text into structured graphs that can be used by downstream NLP tasks.

% These benefits have motivated a growing body of work on automatic text-to-KG extraction, which converts unstructured text into structured graphs that can be used by downstream NLP tasks \citep{kg-llm-survey-arxiv}.

Several LLM-based methods have been proposed for text-to-KG extraction, including KGGen \citep{kggen}, RAKG \citep{rakg}, GraphRAG \citep{graphrag_microsoft}, PiVe  \citep{pive}, and iText2KG \citep{itext2kg}. Despite differences in their specific designs, all of these methods share a common pattern: entities are first extracted from the source text, and relations are then derived only between pairs of those extracted entities. We refer to this shared pattern as the \emph{entity-first} pipeline. 

% To evaluate the effectiveness of these methods, \citep{gbertscore} has proposed structual similarity based GBertScore metrics which is an extension of the \citep{bertscore}. \citep{graphjudge} later argued that these metrics are not suitable of text to kg scenarios, and \citep{kggen} proposed the MINE score which quantify the knowledge coverage. 

The \emph{entity-first} design favors precision because every relation is grounded against an explicit entity set, but it carries a structural consequence. Any entity missed in the first stage is lost entirely, since no relation involving it can be recovered later. Equally, any fact in source text that describes a single entity in isolation, without a second entity to pair with, falls outside the relation-extraction step by construction. Together, these two failure modes produce a systematic coverage gap, particularly on facts that involve a single entity. To address this gap, we propose \textbf{HiReKG}, a high-recall text-to-KG extraction framework that reduces reliance on upfront entity extraction, allowing the LLM to surface facts about each entity directly from the source text, including those that do not pair with a second entity. However, removing the entity-pair constraint can leave the graph disconnected and repetitive, with the same entity scattered across multiple unlinked mentions. HiReKG resolves this with an object resolution module that restores connectivity by linking these mentions into shared graph nodes.

% However, Dropping the entity-pair constraint introduces a new difficulty: when an object is a free-form phrase rather than a known entity, any other entities it mentions remain trapped inside that phrase and never become graph nodes, leaving the resulting graph disconnected. HiReKG resolves this with an object resolution module that detects such entities inside object phrases and rewrites the affected triples so that every mentioned entity appears as its own node, restoring graph connectivity.

\begin{figure*}[t!]
  \centering
  \includegraphics[width=\textwidth]{images/flow.pdf}
  \caption{Comparison of the entity-first paradigm (top) and proposed \textbf{HiReKG} (bottom).}
  \label{fig:pipeline}
\end{figure*}

The main contributions of this paper are as follows:
\begin{itemize}
    \item \textbf{HiReKG}, a high-recall text-to-KG extraction framework that, unlike the dominant entity-first paradigm, recovers single-entity facts while preserving competitive precision.
    \item A single-entity vs.\ multi-entity decomposition showing that entity-first baselines systematically under-recover single-entity facts, a gap that HiReKG largely closes.
\end{itemize}
% We evaluate HiReKG against the two strongest LLM-based baselines, KGGen and RAKG, on three datasets spanning general, scientific, and technical domains, across three language models. HiReKG showed better recall across all configurations with a minor precision trade-off. We also showed a decomposition analysis confirming that this gain is disproportionately driven by the recovery of single-entity facts that entity-first pipelines structurally miss.

% \begin{figure*}[t!]
%   \centering
%   \includegraphics[width=\textwidth]{images/flow.png}
%   \caption{Comparison of the entity-first paradigm (top) and proposed \textbf{HiReKG} (bottom).}
%   \label{fig:pipeline}
% \end{figure*}

% The main contributions of this paper are as follows:
% \begin{itemize}
%     \item A text-to-KG extraction framework that, unlike the dominant entity-first paradigm, recovers facts describing a single entity in isolation while preserving competitive precision.
%     \item A decomposition analysis that attributes recall gains to specific fact categories, providing the evidence that entity-first pipelines structurally underperform on single-entity information.
%     \item Two complementary evaluation metrics adapted to the text-to-KG setting: a Faithfulness Score and a Unique Information Ratio, that fill the gap left by evaluation studies in prior work.
% \end{itemize}

% \section{Introduction}

% Knowledge graphs (KGs) are a structured representation of information that stores entities and the relations between them, enabling strong multi-hop connectivity that unstructured text lacks. Incorporating KGs into RAG, recommendation systems, text summarization, and other NLP tasks has been shown to substantially improve their ability to capture scattered context from source data (\citealp{llm-kg-emnlp-survey}, \citealp{recommendation_kg}, \citealp[]{text_summarization_kg}). KGs have also been used to mitigate hallucinations in LLM responses by grounding generation in verifiable facts (\citealp[]{hallucination-kg-survey}). These benefits have motivated a growing body of work on automatic text-to-KG extraction, which converts unstructured text into structured graphs that can be consumed by downstream NLP tasks \citep{kg-llm-survey-arxiv}.

% Three properties largely determine how useful a KG is for downstream NLP tasks: high knowledge coverage, low redundancy, and robustness across domains [CITE: KG survey, e.g., Hogan et al. or Pan et al. roadmap]. Each generation of text-to-KG approaches has attempted to satisfy these properties to varying degrees. Early approaches relied on predefined ontologies and schemas [CITE: Chepurova et al., Wikontic]; these produced compact, low-redundancy graphs but were tied to the chosen ontology and did not generalize across domains. Later approaches introduced learned components, such as trained entity extractors or triple classifiers [CITE: Anuyah et al.] [CITE: Huang et al., GraphJudger], which improved over rule-based pipelines but inherited the domain limitations of their training data and continued to lose information on different domains. LLM-based approaches leverage the broad world knowledge and language understanding of modern LLMs and have therefore become the dominant paradigm, largely closing the robustness gap [CITE: cluster — KGGen, RAKG, GraphRAG, or a recent LLM-for-KG survey]. However, the most prominent LLM-based methods still fall short on the remaining two properties: knowledge coverage and redundancy.

% Several LLM-based methods have been proposed for text-to-KG extraction, including KGGen [CITE: Mo et al.], RAKG [CITE: Zhang et al.], GraphRAG [CITE: Microsoft Research], and IT2KG [CITE]. Despite differences in their specific designs, all of these methods share a common underlying pipeline: entities are first extracted from the source text, and relations are then derived only between pairs of those extracted entities (Fig. 1). We refer to this shared pattern as the entity-first pipeline. Among these methods, KGGen and RAKG have been shown to outperform the others on standard text-to-KG benchmarks [CITE: Mo et al.] [CITE: Zhang et al.], and we therefore use them as the strongest representatives of the entity-first paradigm for our comparison. KGGen prompts an LLM to extract entities and relations from text and then applies an iterative clustering step that merges semantically equivalent entities and relations into canonical forms, producing compact, low-redundancy graphs with high precision. RAKG operates at the document level by first extracting pre-entities from text chunks and then using each pre-entity as a query to retrieve all relevant chunks from the source document, so that the LLM extracts relations with full document-level context rather than from isolated chunks, mitigating long-context forgetting and coreference errors. The entity-first design is precision-friendly, but it carries a structural consequence — any fact about an entity that does not pair with another extracted entity, or any entity missed in the first stage, cannot be recovered downstream. The result is a systematic coverage gap on facts that describe a single entity in isolation. RAKG additionally exhibits a redundancy problem, generating roughly five to seven times more triples than other approaches, a large fraction of which are paraphrased or duplicated, inflating the graph and degrading its usefulness for downstream tasks. These two observations — the shared coverage gap and the redundancy issue — directly motivate the design of HiReKG.

% In this paper, we propose \textbf{HiReKG}, a high-recall text-to-KG extraction framework that addresses the coverage gap of the entity-first paradigm. Instead of binding relation extraction to a pre-extracted entity set, HiReKG reduces the dependence on the upfront entity extraction step, allowing the LLM to surface information about each entity directly from the source text. To keep the resulting triples cohesive and free of fragmentation that would otherwise arise from removing the entity-set anchor, HiReKG introduces an object resolution module that canonicalizes triples after extraction. We evaluate HiReKG on three datasets that span general, scientific, and technical domains (MINE-1, SciERC, and Re-DocRED), across three language models (GPT-4o, Qwen3-14B, and Qwen3-8B), using two existing metrics (MINE Score and G-BERTScore) and two newly proposed metrics (Faithfulness Score and Unique Information Ratio). Across all configurations, HiReKG improves recall over the strongest baseline by 6--48\% with a precision trade-off of only 1.5--8\%, and produces graphs that are structurally comparable to the baselines while exhibiting lower redundancy.

% The main contributions of this paper are as follows:
% \begin{itemize}
%     \item We propose HiReKG, a text-to-KG extraction pipeline that reduces reliance on upfront entity extraction and incorporates an object resolution module to produce cohesive, high-coverage knowledge graphs.
%     \item We introduce two new evaluation metrics --- the Faithfulness Score for measuring whether predicted triples are grounded in the source text, and the Unique Information Ratio for quantifying redundancy in the extracted graph.
%     \item We provide a comprehensive empirical comparison of HiReKG against KGGen and RAKG across three datasets and three LLM backbones, demonstrating consistent gains in recall while preserving competitive precision and structural quality.
% \end{itemize}

% The remainder of this paper is organized as follows. Section~\ref{sec:related-work} reviews related work in text-to-KG extraction. Section~\ref{sec:method} describes the HiReKG pipeline in detail. Section~\ref{sec:experimental-setup} presents the experimental setup, including datasets, evaluation metrics, and language models. Section~\ref{sec:results} reports and analyzes the results. Section~\ref{sec:conclusion} concludes and discusses directions for future work.


% \section{Related Work}
% \label{sec:related-work}

% Text-to-KG extraction methods can be broadly grouped into two segments: semi-automatic approaches that rely on external structure such as an ontology or a trained module, and fully automatic LLM-based approaches that construct the KG end-to-end. Following the scope established in Section~\ref{sec:introduction}, our quantitative comparison is restricted to LLM-based methods, but we briefly review both segments to position HiReKG within the broader literature.

% \subsection{Semi-Automatic Approaches}

% Early text-to-KG methods depend on an external source of structure to constrain the extraction process. \citet{angeli2015openie} introduce OpenIE, which uses linguistic patterns over dependency parses to extract subject--predicate--object triples directly from text without requiring a schema, but at the cost of noisy and inconsistent triples that are difficult to consume downstream. \citet{chepurova2026wikontic} propose Wikontic, which aligns extracted entities and relations with the Wikidata ontology, producing schema-consistent graphs that are easy to query but narrow in coverage — any concept outside the ontology is dropped. \citet{anuyah2025automated} use trained models for sentence complexity classification and triplet filtering, and \citet{huang2025graphjudger} fine-tune an LLM to act as a graph judge that validates candidate triples before they are added to the graph. Both improve over rule-based pipelines on in-domain data but inherit the domain limitations of their training corpora and degrade when applied to text outside the training distribution.

% Across this segment, the common limitation is a lack of robustness across domains: ontology-aligned methods are bound to a fixed schema, and methods with trained components are bound to their training distribution. Neither generalizes cleanly to arbitrary input text, which is the setting our work targets.

% \subsection{LLM-Based Approaches}

% LLM-based methods exploit the broad world knowledge and language understanding of large pre-trained models to extract entities and relations without external schemas or task-specific training. \citet{microsoft2024graphrag} introduce GraphRAG, which constructs a hierarchical knowledge graph from text chunks and uses it to support retrieval-augmented generation. \citet{mo2026kggen} propose KGGen, which prompts an LLM to extract entities and relations and then applies an iterative LLM-based clustering step to merge semantically equivalent entities and relations into canonical forms, producing compact graphs with high precision. The authors also release the MINE benchmark, which measures how much of the source-text information is recoverable from the extracted graph. \citet{zhang2025rakg} propose RAKG, a document-level construction framework that first extracts pre-entities from text chunks and then uses each pre-entity as a query to retrieve all relevant chunks from the document, so that the LLM extracts relations with full document-level context rather than from a single chunk. This design mitigates long-context forgetting and coreference errors that arise when LLMs process documents chunk by chunk.

% Despite their different designs, all of these methods follow the entity-first pipeline described in Section~\ref{sec:introduction}: entities are extracted upfront and relations are derived only between those entities. This shared structural choice is the source of the coverage gap our work addresses.

% \subsection{Excluded Baselines}
% \label{sec:excluded-baselines}

% We restrict our quantitative comparison to KGGen and RAKG as the strongest representatives of the LLM-based, entity-first paradigm. Prior work has established that these two methods outperform other LLM-based extractors on standard text-to-KG benchmarks: RAKG reports a 6.2-point improvement over GraphRAG on the MINE benchmark \citep{zhang2025rakg}, and KGGen reports comparable advantages on its own benchmark against earlier LLM-based extractors \citep{mo2026kggen}. We therefore treat GraphRAG and similar methods such as IT2KG \citep{itk2g-citation} as part of the broader landscape that we contextualize but do not benchmark directly. We note that GraphRAG is widely deployed in downstream RAG pipelines; comparing HiReKG to GraphRAG in an end-to-end RAG setting, rather than on KG-quality metrics alone, is left to future work.

% \subsection{Evaluation of Text-to-KG Methods}

% Evaluation of extracted KGs has historically lacked a consistent recall measure across datasets. Graph-level structural metrics such as G-BERTScore \citep{saha2021explagraphs} compute precision and recall via maximum-weight bipartite matching between predicted and reference triples, but these measure structural overlap with a reference graph rather than how much of the source-text information is captured. The MINE benchmark \citep{mo2026kggen} addresses this gap by checking whether a set of manually verified facts can be inferred from the extracted graph, but it is tied to a single dataset. Faithfulness — whether the triples in the extracted graph are actually supported by the source text — has been treated only implicitly in prior text-to-KG work, although the broader RAG and summarization literature has developed faithfulness evaluation frameworks based on LLM-as-judge protocols \citep[CITE: faithfulness-eval-survey-or-paper]. Redundancy in the extracted graph has similarly received little attention, despite its direct impact on downstream graph-based RAG pipelines. The Faithfulness Score and Unique Information Ratio we introduce in Section~\ref{sec:experimental-setup} are designed to fill these two gaps and to apply consistently across datasets.




\section{HiReKG}
\label{sec:method}
HiReKG extracts a knowledge graph from plain text in three stages: subject extraction, property extraction, and object resolution, followed by the refinement stage of \citet{kggen}. The source document is first split into chunks, and each stage runs over all chunks before the next begins. Figure~\ref{fig:pipeline} contrasts HiReKG with the entity-first paradigm: subject extraction mirrors the entity-extraction step of prior pipelines, but HiReKG then decouples property extraction from the entity-pair anchor and adds an object resolution step that restores graph connectivity before refinement. A full end-to-end run on a sample paragraph is provided in Appendix~\ref{app:walkthrough}.

\subsection{Subject Extraction}
\label{sec:method-subject}
Following prior LLM-based extractors \citep{kggen, rakg}, we prompt an LLM on each chunk to identify the set of entities mentioned in the source document. Each entity in this set serves as a \emph{subject} around which information is later organized.


\subsection{Property Extraction}
\label{sec:method-property}
For each chunk, the LLM is prompted to gather all information about the set of subjects identified in Section~\ref{sec:method-subject}. The output is a set of \emph{triples} \texttt{(subject, relation, object)}, where the subject is one of the extracted entities, the relation is a free-form predicate, and the object is a natural-language phrase describing a property, attribute, or related entity. Crucially, and unlike entity-first methods, the object is not constrained to be another extracted entity: it may be a date, a quantity, a descriptive clause, or a phrase that itself references one or more entities.

This relaxation is what lets HiReKG recover single-entity facts. A statement such as ``The Eiffel Tower is 330 metres tall'' yields the triple \texttt{(Eiffel Tower, height, 330 metres)} keyed on a single subject, with no second entity required. 

The trade-off is structural. When an object is a phrase that references other entities, those entities remain trapped inside the object string rather than appearing as proper graph nodes; for instance, ``Marie Curie won the Nobel Prize in Physics and the Nobel Prize in Chemistry'' is extracted as a single triple whose object embeds two prize entities that should each be a node in their own right. As a result, the graph initially fragments into per-subject \emph{stars}, as shown in the left pane of Figure~\ref{fig:object-resolution}. Object resolution repairs this fragmentation.

% \subsection{Property Extraction}
% \label{sec:method-property}

% For each text chunk, the LLM is prompted to gather, for every subject identified in Section~\ref{sec:method-subject}, all information about that subject that is supported by the chunk. The output is a set of triples in which the subject is one of the extracted entities and the object is a natural-language statement describing a property, attribute, or relation. Unlike entity-first methods, the object is not constrained to be another extracted entity; it can be a free-form phrase such as a date, a quantity, a descriptive clause, or a phrase that references a different entity.

% This design is the primary mechanism by which HiReKG reduces the entity-first coverage gap. A fact such as ``Einstein won the 1921 Nobel Prize in Physics'' is recovered as a triple keyed on \emph{Einstein}, regardless of whether ``Nobel Prize'' was identified as an entity in the first stage. The trade-off is structural: because objects are not entity-aligned by construction, the resulting triples can contain entities embedded inside descriptive object strings, and the graph initially fragments into per-subject ``stars'' rather than a connected structure, as shown in the left pane of the Figure \ref{fig:object-resolution}. Object resolution, described next, addresses this fragmentation.

\begin{figure}[t]
  \centering
  \includegraphics[width=0.75\columnwidth]{images/obj_res.pdf}
  \caption{Object resolution links per-subject triples (left) by merging coreferent objects into shared nodes (right).}
  \label{fig:object-resolution}
\end{figure}




\subsection{Object Resolution}
\label{sec:method-resolution}
Object resolution rewrites raw triples so that any extracted entity appearing inside an object string becomes a proper object (Figure~\ref{fig:object-resolution}). For each raw triple, the object is normalized (lowercased, lemmatized, with determiners stripped) and matched against the extracted entity set using word-boundary string matching; triples with no match are forwarded unchanged. Triples with one embedded entity can be split apart by a fixed rule, while triples with multiple embedded entities can be split in several different ways and need the LLM to choose. We therefore use two repair procedures.
\paragraph{Single embedded entity.} The object is tokenized and decomposed into a \emph{prefix} (text before the entity), the entity span, and a \emph{suffix} (text after the entity). In the common case, only a prefix is present: the prefix is folded into the relation and the entity becomes the new object, producing $(s, r', e)$ where $r'$ is the relation with the prefix appended. This deterministic path handles the majority of single-entity cases without an LLM call. When a suffix is present, an LLM rewrites the triple so that the entity becomes the object and any remaining content is folded into the relation.
\paragraph{Multiple embedded entities.} The LLM is prompted with the original triple, the list of detected entities, and a structured analysis of each entity's surrounding context. It selects one of three decomposition strategies and emits a set of triples in which every embedded entity appears as a proper object. \textsc{Hub} handles parallel lists, such as the two-Nobel-Prize triple from Section~\ref{sec:method-property}; \textsc{Chain} handles sequential connectives; and \textsc{Nested} handles hierarchical phrasing.

All LLM-rewritten triples are validated for entity coverage and field completeness, reverting to the original on failure. By promoting embedded entities to graph nodes, object resolution converts the per-subject \emph{stars} of Section~\ref{sec:method-property} into a connected graph without imposing the entity-first constraint at extraction time. Worked examples for each pattern (single-entity, \textsc{Hub}, \textsc{Chain}, \textsc{Nested}) appear in Appendix~\ref{app:obj-res-examples}.

% All LLM-rewritten triples are validated to ensure that detected entities are surfaced and that no field is empty; triples that fail validation revert to the original. By promoting embedded entities to graph nodes, object resolution converts the per-subject \emph{stars} of Section~\ref{sec:method-property} into a connected graph, without imposing the entity-first constraint at extraction time. Worked examples for each pattern (single-entity, \textsc{Hub}, \textsc{Chain}, and \textsc{Nested}) are provided in Appendix~\ref{app:obj-res-examples}.


% \subsection{Object Resolution}
% \label{sec:method-resolution}

% Object resolution rewrites raw triples so that any extracted entity appearing inside an object string becomes a proper object in its own right (Figure~\ref{fig:object-resolution}). For each raw triple, the object is normalized (lowercased, lemmatized, with determiners stripped) and checked against the extracted entity set using word-boundary matching. Triples whose objects contain no extracted entity are forwarded unchanged. Triples whose objects contain one or more embedded entities are routed to one of two repair procedures depending on the number of entities found.

% \paragraph{Single embedded entity.} When exactly one entity is embedded, the object is tokenized and decomposed into a \emph{prefix} (text before the entity), the entity span, and a \emph{suffix} (text after the entity). If only a prefix is present, the prefix is folded into the relation and the entity becomes the new object, producing a triple of the form $(s, r{+}\mathrm{prefix}, e)$. If a suffix is present, an LLM call rewrites the triple so that the entity becomes the object and any remaining content is redistributed into the relation. The deterministic prefix-only path handles the majority of single-entity cases without an LLM call.

% \paragraph{Multiple embedded entities.} When two or more entities are embedded in the object, an LLM is prompted with the original triple, the list of detected entities, and a structured analysis of each entity's surrounding context. The LLM selects a decomposition strategy --- \textsc{Hub} for parallel lists, \textsc{Chain} for sequential connectives, or \textsc{Nested} for hierarchical phrasing --- and emits a set of triples in which every embedded entity appears as a proper object. The output is validated to ensure that all detected entities are surfaced and that no triple has an empty field; triples that fail validation revert to the original.

% Object resolution is the enabling step that makes the property-extraction design viable: by promoting embedded entities to graph nodes, it converts the disconnected per-subject stars produced by Section~\ref{sec:method-property} into a connected graph without re-introducing the entity-first constraint at extraction time.

% \subsection{Refinement}
% \label{sec:method-refinement}
% The resolved triples are passed through the iterative LLM-based refinement of KGGen \citep{kggen}, which merges semantically equivalent entities into canonical forms, using embedding models, yielding the final knowledge graph.

\section{Experiments}
\label{sec:experiments}


\begin{table*}[t!]
\centering
\small
\setlength{\tabcolsep}{3pt}
\begin{tabular}{ll|cccc|cccc|cccc}
\toprule
& & \multicolumn{4}{c|}{\textbf{MINE-1}} & \multicolumn{4}{c|}{\textbf{SciERC}} & \multicolumn{4}{c}{\textbf{Re-DocRED}} \\
\textbf{LLM} & \textbf{Method} & Single & Multi & Gap & Overall & Single & Multi & Gap & Overall & Single & Multi & Gap & Overall \\
\midrule
\multirow{3}{*}{GPT-4o} 
& KGGen      & 59.44 & 66.07 & $-6.63$ & 65.05 & 62.74 & 68.36 & $-5.62$  & 66.17 & 35.71 & 53.89 & $-18.18$ & 49.33 \\
& RAKG       & 63.08 & 77.68 & $-14.60$ & 74.33 & 35.71 & 35.29 & $+0.42$  & 35.46 & 37.50 & 56.29 & $-18.79$ & 51.57 \\
& \textbf{HiReKG} & \textbf{91.28} & \textbf{94.03} & $-2.75$  & \textbf{93.40} & \textbf{81.03} & \textbf{82.02} & $-0.99$  & \textbf{81.65} & \textbf{80.36} & \textbf{80.24} & $+0.12$  & \textbf{80.27} \\
\midrule
\multirow{3}{*}{Qwen3-14B} 
& KGGen      & 45.93 & 62.72 & $-16.79$ & 58.87 & 49.91 & 56.45 & $-6.54$  & 53.91 & 50.00 & 45.51 & $+4.49$  & 46.64 \\
& RAKG       & 55.81 & 62.98 & $-7.17$  & 61.33 & 43.50 & 46.86 & $-3.36$  & 45.55 & 37.50 & 60.48 & $-22.98$ & 54.71 \\
& \textbf{HiReKG} & \textbf{88.08} & \textbf{91.70} & $-3.62$  & \textbf{90.87} & \textbf{76.95} & \textbf{79.93} & $-2.98$  & \textbf{78.77} & \textbf{82.14} & \textbf{79.04} & $+3.10$  & \textbf{79.82} \\
\midrule
\multirow{3}{*}{Qwen3-8B} 
& KGGen      & 46.29 & 58.69 & $-12.40$ & 55.85 & 53.73 & 51.93 & $+1.80$  & 52.63 & 42.86 & 49.10 & $-6.24$  & 47.53 \\
& RAKG       & 62.50 & 77.16 & $-14.66$ & 73.80 & 64.64 & 63.62 & $+1.02$  & 64.02 & 41.07 & 59.88 & $-18.81$ & 55.16 \\
& \textbf{HiReKG} & \textbf{77.03} & \textbf{85.29} & $-8.26$  & \textbf{83.40} & \textbf{66.90} & \textbf{68.91} & $-2.01$  & \textbf{68.13} & \textbf{60.71} & \textbf{65.27} & $-4.56$  & \textbf{64.13} \\
\bottomrule
\end{tabular}
\caption{MINE Score (recall, \%) decomposed by entity count. \textbf{Single}/\textbf{Multi}: facts mentioning one or multiple entities. \textbf{Gap} = Single $-$ Multi. \textbf{Overall}: full-dataset score. \textbf{Bold} marks the best per cell.}
\label{tab:recall}
\end{table*}

% \begin{table*}[t!]
% \centering
% \small
% \setlength{\tabcolsep}{4pt}
% \begin{tabular}{lll|ccc|ccc|ccc}
% \toprule
% & & & \multicolumn{3}{c|}{\textbf{MINE-1}} & \multicolumn{3}{c|}{\textbf{SciERC}} & \multicolumn{3}{c}{\textbf{Re-DocRED}} \\
% \textbf{LLM} & \textbf{Method} & & Single & Multi & Gap & Single & Multi & Gap & Single & Multi & Gap \\
% \midrule
% \multirow{3}{*}{GPT-4o} 
% & KGGen      & & 48.25   & 66.34   & $-18.09$   & 62.74 & 68.36 & $-5.62$  & 35.71 & 53.89 & $-18.18$ \\
% & RAKG       & & 63.08 & 77.68 & $-14.60$ & 35.71 & 35.29 & $+0.42$  & 37.50 & 56.29 & $-18.79$ \\
% & \textbf{HiReKG} & & \textbf{91.28} & \textbf{94.03} & $-2.75$  & \textbf{81.03} & \textbf{82.02} & $-0.99$  & \textbf{80.36} & \textbf{80.24} & $+0.12$  \\
% \midrule
% \multirow{3}{*}{Qwen3-14B} 
% & KGGen      & & 45.93 & 62.72 & $-16.79$ & 49.91 & 56.45 & $-6.54$  & 50.00 & 45.51 & $+4.49$  \\
% & RAKG       & & 55.81 & 62.98 & $-7.17$  & 43.50 & 46.86 & $-3.36$  & 37.50 & 60.48 & $-22.98$ \\
% & \textbf{HiReKG} & & \textbf{88.08} & \textbf{91.70} & $-3.62$  & \textbf{76.95} & \textbf{79.93} & $-2.98$  & \textbf{82.14} & \textbf{79.04} & $+3.10$  \\
% \midrule
% \multirow{3}{*}{Qwen3-8B} 
% & KGGen      & & 46.29 & 58.69 & $-12.40$ & 53.73 & 51.93 & $+1.80$  & 42.86 & 49.10 & $-6.24$  \\
% & RAKG       & & 62.50 & 77.16 & $-14.66$ & 64.64 & 63.62 & $+1.02$  & 41.07 & 59.88 & $-18.81$ \\
% & \textbf{HiReKG} & & \textbf{77.03} & \textbf{85.29} & $-8.26$  & \textbf{66.90} & \textbf{68.91} & $-2.01$  & \textbf{60.71} & \textbf{65.27} & $-4.56$  \\
% \bottomrule
% \end{tabular}
% \caption{Single-entity vs. multi-entity decomposition of MINE Score (\%). Entity-first baselines (KGGen, RAKG) show a systematic single-entity deficit; HiReKG largely closes the gap and recovers the absolute scores by a wide margin. \textbf{Bold} marks the best method per cell.}
% \label{tab:single-entity}
% \end{table*}


\begin{table}[t!]
\centering
\small
\setlength{\tabcolsep}{4pt}
\begin{tabular}{ll|ccc}
\toprule
\textbf{LLM} & \textbf{Method} & \textbf{MINE-1} & \textbf{SciERC} & \textbf{Re-DocRED} \\
\midrule
\multirow{3}{*}{GPT-4o} 
& KGGen           & \textbf{90.23} & \textbf{83.44} & \textbf{85.05} \\
& RAKG            & 57.38          & 61.11          & 64.01 \\
& \textbf{HiReKG} & 81.06          & 78.77          & 77.46 \\
\midrule
\multirow{3}{*}{Qwen3-14B} 
& KGGen           & \textbf{93.96} & \textbf{93.73} & \textbf{95.70} \\
& RAKG            & 81.38          & 78.01          & 76.24 \\
& \textbf{HiReKG} & 90.29          & 87.17          & 91.82 \\
\midrule
\multirow{3}{*}{Qwen3-8B} 
& KGGen           & \textbf{91.83} & \textbf{90.38} & \textbf{94.74} \\
& RAKG            & 80.90          & 73.99          & 73.79 \\
& \textbf{HiReKG} & 89.72          & 83.12          & 88.12 \\
\bottomrule
\end{tabular}
\caption{Faithfulness Score (precision, \%). KGGen leads precision in every cell; HiReKG trails by 2--9 pp while delivering the largest recall gains (Table~\ref{tab:recall}). \textbf{Bold} marks the best per cell.}
% \label{tab:faith}
\label{tab:faith}
\end{table}

% \begin{table*}[t!]
% \centering
% \small
% \setlength{\tabcolsep}{4pt}
% \begin{tabular}{lll|cc|cc|cc}
% \toprule
% & & & \multicolumn{2}{c|}{\textbf{MINE-1}} & \multicolumn{2}{c|}{\textbf{SciERC}} & \multicolumn{2}{c}{\textbf{Re-DocRED}} \\
% \textbf{LLM} & \textbf{Method} & & MINE & Faith. & MINE & Faith. & MINE & Faith. \\
% \midrule
% \multirow{3}{*}{GPT-4o} 
% & KGGen      & & 63.33   &  \textbf{90.23}   & 66.17 & \textbf{83.44} & 49.33 & \textbf{85.05} \\
% & RAKG       & & 74.33 & 57.38 & 35.46 & 61.11 & 51.57 & 64.01 \\
% & \textbf{HiReKG} & & \textbf{93.40} & 81.06 & \textbf{80.19} & 78.77 & \textbf{80.27} & 77.46 \\
% \midrule
% \multirow{3}{*}{Qwen3-14B} 
% & KGGen      & & 58.87 & \textbf{93.96} & 53.91 & \textbf{93.73} & 46.64 & \textbf{95.70} \\
% & RAKG       & & 61.33 & 81.38 & 45.55 & 78.01 & 54.71 & 76.24 \\
% & \textbf{HiReKG} & & \textbf{90.87} & 90.10 & \textbf{78.77} & 87.17 & \textbf{79.82} & 91.82 \\
% \midrule
% \multirow{3}{*}{Qwen3-8B} 
% & KGGen      & & 55.85 & \textbf{91.83} & 52.63 & \textbf{90.38} & 47.53 & \textbf{94.74} \\
% & RAKG       & & 73.80 & 80.90 & 64.02 & 73.99 & 55.16 & 73.79 \\
% & \textbf{HiReKG} & & \textbf{83.40} & 89.72 & \textbf{68.13} & 83.12 & \textbf{64.13} & 88.12 \\
% \bottomrule
% \end{tabular}
% \caption{Main results across nine LLM--dataset configurations. \textbf{MINE} = MINE Score (recall, \%). \textbf{Faith.} = Faithfulness Score (precision, \%). HiReKG achieves the highest MINE Score in every configuration. \textbf{Bold} marks the best method per cell.}
% \label{tab:main}
% \end{table*}


\subsection{Datasets}
\label{sec:datasets}
We evaluate on three datasets spanning distinct text styles. \textbf{MINE-1} \citep{kggen} provides 100 LLM-generated documents paired with manually verified atomic facts that serve as a recall ground truth. \textbf{SciERC} \citep{scierc} contains 100 scientific abstracts    annotated with entities and relations, representing dense technical prose. \textbf{Re-DocRED} \citep{redocred} is a cleaned version of DocRED \citep{docred} built from Wikipedia documents; we use the 15-document subset introduced by \citet{lkd-kgc}. All three datasets have been used in recent text-to-KG work \citep{graphjudge, rakg}, enabling direct comparison with prior results.

% \subsection{Datasets}
% \label{sec:datasets}

% We evaluate on three datasets spanning distinct text styles. \textbf{MINE-1} \citep{kggen} consists of 100 LLM-generated documents paired with manually verified atomic facts that serve as a recall ground truth. \textbf{SciERC} \citep{scierc} contains 100 scientific abstracts annotated with entities and relations, representing dense technical prose. \textbf{Re-DocRED} \citep{redocred} is a revised, cleaned version of DocRED \citep{docred}, built from Wikipedia documents with corrected relation annotations; we use a 15-document subset given its longer document length. Together, the three datasets cover general, scientific, and encyclopedic domains and let us check whether the observed trends hold across text styles. These three datasets have also been used by recent text-to-KG work \citep{graphjudge, rakg, lkd-kgc}, enabling direct comparison with prior results.


\subsection{Baselines and Metrics}
\label{sec:baselines}
\paragraph{Baselines.} We compare HiReKG against \textbf{KGGen} \citep{kggen} and \textbf{RAKG} \citep{rakg}, the two strongest representatives of the entity-first paradigm.
\paragraph{Metrics.} We report recall and precision via two LLM-as-judge metrics. \textbf{MINE Score} \citep{kggen} (recall) is the ratio of ground truth atomic facts that can be inferred from the extracted graph. \textbf{Faithfulness Score} (precision), adapted from RAG evaluation \citep{ragas}, measures the fraction of extracted triples supported by the source text. As a supporting measure, the deterministic \textbf{Unique Information Ratio (UIR)} captures redundancy as the fraction of non-overlapping triples; UIR results are reported in Appendix~\ref{app:uir}. Full prompts and configurations of all metrics are in Appendix~\ref{app:metrics}.
\paragraph{Models.} We run each method with three LLM backbones: \textbf{GPT-4o} \citep{gpt4o}, \textbf{Qwen3-14B} \citep{qwen3_14b}, and \textbf{Qwen3-8B} \citep{qwen3_8b}, spanning one flagship closed model and two open-weight models of different scales. Exact versions and hyperparameters are listed in Appendix~\ref{app:hardware}.

% \subsection{Baselines and Metrics}
% \label{sec:baselines}

% \paragraph{Baselines.} We compare HiReKG against \textbf{KGGen} \citep{kggen} and \textbf{RAKG} \citep{rakg}, the two strongest representatives of the entity-first paradigm.

% \paragraph{Metrics.} We report three metrics. \textbf{MINE Score} \citep{kggen} measures recall by checking whether each atomic fact in the dataset can be inferred from the extracted graph, scored by an LLM judge. \textbf{Faithfulness Score} measures precision as the fraction of extracted triples that are entailed by the source text, also scored by an LLM judge. \textbf{Unique Information Ratio (UIR)} measures redundancy as the fraction of triples in the graph that carry non-overlapping information. Faithfulness and UIR are adapted from the broader LLM-as-judge evaluation literature \citep{llm-kg-emnlp-survey} to the text-to-KG setting; full description and prompt templates are provided in Appendix~\ref{app:metrics}.

% \paragraph{Models.} We run each method with three LLM backbones: \textbf{GPT-4o} \citep{gpt4o}, \textbf{Qwen3-14B} \citep{qwen3_14b}, and \textbf{Qwen3-8B} \citep{qwen3_8b}. This combination spans one flagship closed model and two open-weight models of different scales, letting us check whether the observed trends are model-dependent. Hyperparameters and exact model versions are listed in Appendix~\ref{app:hardware}.


\subsection{Single-Entity Decomposition Results}
\label{sec:single-entity}
Our central hypothesis is that entity-first pipelines structurally under-recover facts that describe a single entity in isolation. To test this directly, we classify each atomic fact in the ground truth as \emph{single-entity} (mentioning one of the document's entities) or \emph{multi-entity} (connecting two or more), and report the MINE Score on each subset.\footnote{Classification uses ground-truth entity annotations on SciERC and Re-DocRED, aand GPT-5~\citep{gpt5} on MINE-1 with a spaCy NER sanity check (76.8\% agreement); see Appendix~\ref{app:datasets}.}

Table~\ref{tab:recall} reports the decomposition, and the pattern is consistent with the hypothesis. KGGen and RAKG show a systematic single-entity deficit: their single-entity score is lower than their multi-entity score in 14 of 18 configurations, with average gaps of $-7.4$ pp and $-11.0$ pp respectively. The deficit is sharpest on Re-DocRED, where both baselines collapse to single-entity scores of $35.7\%$ and $37.5\%$ with GPT-4o while remaining above $53\%$ on multi-entity facts. HiReKG, in contrast, is near-symmetric (average gap $-2.4$ pp), and its absolute single-entity score exceeds the strongest baseline by $+24.1$ pp on average (up to $+42.9$ pp), with the largest gains precisely where the baselines collapse. The recall advantage is therefore not a generic uplift: it is concentrated on the class of facts the entity-first design undercaptures.

% \subsection{Single-Entity Decomposition}
% \label{sec:single-entity}

% Our central hypothesis is that entity-first pipelines structurally under-recover facts that describe a single entity in isolation. To test this directly, we classify each atomic fact in the recall ground truth as \emph{single-entity} (the fact mentions only one of the document's entities) or \emph{multi-entity} (the fact connects two or more entities), and report the MINE Score separately on the two subsets. Classification uses each dataset's ground truth entity annotations where available (SciERC, Re-DocRED) and GPT-5~\citep{gpt5} validated against spaCy NER otherwise (MINE-1) (Appendix~\ref{app:datasets}).

% Table~\ref{tab:single-entity} reports the decomposition. The pattern is consistent with the hypothesis. KGGen and RAKG show a systematic deficit on single-entity facts: their single-entity score is lower than their multi-entity score in 14 of 18 configurations, with an average gap of $-8.6$ percentage points (pp) for KGGen and $-11.0$ pp for RAKG. The deficit is sharpest on Re-DocRED, where both baselines collapse to single-entity scores of $35.7\%$ and $37.5\%$ with GPT-4o while still scoring above $53\%$ on multi-entity facts. HiReKG, in contrast, exhibits a near-symmetric profile: the average single-vs-multi gap is $-2.4$ pp, and HiReKG's absolute single-entity score exceeds the strongest baseline by $+23.9$ pp on average (up to $+42.9$ pp) across the nine LLM--dataset configurations, with the largest gains precisely where the baselines collapse. This is direct mechanistic evidence that the recall advantage of HiReKG is not a generic uplift; it is concentrated on the fact category that the entity-first design structurally cannot recover.

\subsection{MINE and Faithfulness Scores}
\label{sec:main-results}
Tables~\ref{tab:recall} and~\ref{tab:faith} report overall recall and precision across the nine LLM--dataset configurations. HiReKG improves recall over the strongest baseline in every configuration, by $+4.1$ to $+29.5$ pp absolute ($+6.4\%$ to $+55.7\%$ relative). The gain is largest on MINE-1 and Re-DocRED with the mid-scale Qwen3-14B, and remains substantial on GPT-4o where the baselines are already at their strongest. RAKG and KGGen alternate as the stronger baseline depending on dataset and model, but HiReKG outperforms both in every cell.

On Faithfulness, KGGen retains the highest score in every configuration, consistent with its strict entity-pair anchoring. HiReKG trails KGGen by $2.1$ to $9.2$ pp, with the gap narrowing on the open-weight models ($2.1$ to $7.3$ pp). HiReKG substantially outperforms RAKG, whose Faithfulness falls as low as $57.4\%$ on MINE-1 with GPT-4o. HiReKG's precision therefore sits between the two baselines while recovering the bulk of KGGen's recall gap, a net improvement for downstream tasks that value coverage.

% \subsection{MINE and Faithfulness Scores}
% \label{sec:main-results}

% Table~\ref{tab:main} reports overall MINE Score and Faithfulness Score across all configurations. HiReKG improves recall over the strongest baseline by $+4.1$ to $+29.5$ pp absolute ($+6.4\%$ to $+55.7\%$ relative), and the gain is positive in every one of the nine configurations. Recall gains are largest on MINE-1 and Re-DocRED with mid-scale models (Qwen3-14B), and remain substantial with GPT-4o where the baselines are already at their strongest. RAKG and KGGen alternate as the stronger baseline depending on dataset and model, but neither approaches HiReKG's coverage in any configuration.

% On Faithfulness, KGGen retains the highest score in all configurations, consistent with its strict entity-pair anchoring. HiReKG trails KGGen by $2.1$ to $7.3$ pp on the open-weight runs while substantially outperforming RAKG, whose Faithfulness drops as low as $57.4\%$ on MINE-1 with GPT-4o due to over-generation of unsupported triples. HiReKG's precision sits between the two baselines, but recovers the bulk of KGGen's recall gap; the net effect on graph quality is a strict improvement for any downstream task that values coverage. On UIR, HiReKG is comparable to KGGen and substantially higher than RAKG, indicating that the recall gains do not come from inflating the graph with duplicate or paraphrased triples (full UIR table in Appendix~\ref{app:uir}). 

% \subsection{Object Resolution Statistics}
% \label{sec:obj-res-stats}

% Section~\ref{sec:method-resolution} described object resolution as the
% \emph{enabling} step that converts the disconnected per-subject stars
% produced by property extraction into a connected graph. We now report
% the operational statistics of this module across all nine LLM--dataset
% configurations to substantiate that claim with hard numbers.

% Table~\ref{tab:obj-res-stats} reports, per configuration, the share of
% raw extractor triples that pass through unchanged (clean), are
% successfully rewritten (repaired), fail validation and revert to their
% original form (failed), and the net change in triple count after
% resolution (yield). Aggregated across the nine configurations, the
% module processes 18{,}521 raw triples and emits 22{,}300 final triples,
% a net yield of $+20.4\%$. Only $27.4\%$ of raw triples are clean on
% arrival; $67.0\%$ require active repair, and $5.6\%$ fail validation and
% are retained as-is. The repair rate is tightly bounded across all nine
% configurations ($57.3$--$74.4\%$), indicating that object resolution
% operates at a stable rate independent of dataset or extractor scale.
% The largest absolute yield ($+38.5\%$ on Re-DocRED with GPT-4o,
% $+37.8\%$ on MINE-1 with GPT-4o) is observed where the flagship
% extractor produces the densest compound objects, which the module
% decomposes into multiple atomic triples; the lowest yield
% ($+9.1\%$ on Re-DocRED with Qwen3-14B) coincides with the highest fail
% rate ($9.5\%$), reflecting the documented fallback path in
% Section~\ref{sec:obj-res} in which validation failures revert to the
% original triple. Across all configurations, the clean rate of
% $16.2$--$42.3\%$ confirms that the module is adaptive rather than
% aggressive: triples that are already well-formed are forwarded
% unchanged, and the graph is not inflated through indiscriminate
% rewriting.

% \begin{table*}[t]
% \centering
% \small
% \setlength{\tabcolsep}{4pt}
% \begin{tabular}{llrrrrrrrrr}
% \toprule
% & & \multicolumn{3}{c}{MINE-1} & \multicolumn{3}{c}{SciERC} & \multicolumn{3}{c}{Re-DocRED} \\
% \cmidrule(lr){3-5}\cmidrule(lr){6-8}\cmidrule(lr){9-11}
% LLM & & Rep.\ \% & Fail \% & Yield \% & Rep.\ \% & Fail \% & Yield \% & Rep.\ \% & Fail \% & Yield \% \\
% \midrule
% \multirow{1}{*}{GPT-4o}    & & 72.0 & 1.5 & $+$37.8 & 57.3 & 0.3 & $+$17.5 & 73.8 & 0.4 & $+$38.5 \\
% \multirow{1}{*}{Qwen3-14B} & & 63.9 & 13.4 & $+$13.3 & 68.7 & 11.7 & $+$14.5 & 74.4 & 9.5 & $+$9.1 \\
% \multirow{1}{*}{Qwen3-8B}  & & 64.9 & 2.4 & $+$11.1 & 68.8 & 2.2 & $+$13.2 & 66.9 & 2.4 & $+$12.2 \\
% \bottomrule
% \end{tabular}
% \caption{Object resolution statistics across nine LLM--dataset
% configurations. Rep.\ \% is the fraction of raw extractor triples that
% the module rewrites; Fail \% is the fraction that fail validation and
% revert to the original; Yield \% is the net change in triple count
% after resolution. The remainder of triples in each cell pass through
% unchanged. Aggregated across configurations, the module repairs
% $67.0\%$ of raw triples, fails on $5.6\%$, and increases the final
% triple count by $20.4\%$.}
% \label{tab:obj-res-stats}
% \end{table*}

\subsection{Redundancy and Cost}
\label{sec:cost-redundancy}
We further measure graph redundancy via UIR and inference cost in total tokens per document. HiReKG achieves the highest overall UIR ($0.71$ vs.\ $0.68$ for KGGen and $0.59$ for RAKG), producing graphs that are less redundant per triple, a property that matters for downstream tasks~\citep{redundanyinkg}. The cost overhead is bounded: HiReKG uses roughly $3\times$ the tokens of KGGen but less than half of RAKG. The additional tokens connect the scattered triples from property extraction, driving HiReKG's recall and UIR gains. Detailed breakdowns are in Appendices~\ref{app:uir},~\ref{app:cost}, and~\ref{app:obj-res}.

\section{Conclusion}
\label{sec:conclusion}
We presented HiReKG, a text-to-KG extraction framework that decouples property extraction from entity-pair anchoring and restores graph connectivity through a dedicated object resolution module. Across nine LLM--dataset configurations, HiReKG achieves the highest recall in every cell while keeping precision within a narrow margin of the strongest entity-first baseline. A single-entity decomposition closes the loop: the systematic deficit on facts that describe a single entity in isolation, long latent in entity-first pipelines, is empirically identified and substantially recovered by HiReKG. The result demonstrates that the coverage ceiling of current text-to-KG systems is not intrinsic to the task but is an artifact of design choices that can be relaxed without sacrificing graph quality.

% \section{Conclusion}
% \label{sec:conclusion}

% We presented HiReKG, a text-to-KG extraction framework that reduces reliance on upfront entity extraction and introduces an object resolution module to keep the resulting triples cohesive. Across three datasets and three LLM backbones, HiReKG improves recall over the strongest entity-first baseline in every configuration while incurring only a small precision trade-off. A single-entity decomposition analysis attributes these gains to the recovery of facts that describe a single entity in isolation, providing direct mechanistic evidence that entity-first pipelines structurally under-recover this class of facts. Future work includes evaluating its impact on downstream graph-based RAG pipelines.

\section*{Limitations}

\paragraph{Downstream evaluation.} We evaluate HiReKG on intrinsic KG-quality metrics (recall, precision, redundancy) rather than on downstream tasks such as graph-based retrieval-augmented generation. A fair downstream comparison would require choosing retrieval methods per domain, which is a separate study in itself.

\paragraph{No multi-seed variance.} The scores in Tables~\ref{tab:recall} and~\ref{tab:faith} are single-run point estimates; multi-seed confidence intervals were not feasible under our compute and API budget. We instead report per-document standard deviations in Appendix~\ref{app:variance}, which capture LLM-judge variability across documents within each configuration.

\paragraph{Single-document scope.} HiReKG produces one KG per input document. Multi-document KG construction is a harder, distinct problem that requires cross-document entity linking and conflict resolution between sources, and is outside the scope of this work.

% \section*{Limitations}

% \paragraph{Downstream evaluation.} We evaluate HiReKG on intrinsic KG-quality metrics rather than on downstream tasks such as graph-based retrieval-augmented generation. A fair downstream comparison would require choosing retrieval methods per domain, which is a separate study in itself.

% \paragraph{No variance estimates.} The results in Tables~\ref{tab:single-entity} and~\ref{tab:main} are single-run point estimates; multi-seed confidence intervals were not feasible under our compute and API budget.

% \paragraph{Single-document scope.} HiReKG produces one KG per input document. Multi-document KG construction is a harder, distinct problem that requires cross-document entity linking and conflict resolution between sources, and is outside the scope of this work.
% \section{Preamble}

% The first line of the file must be
% \begin{quote}
% \begin{verbatim}
% \documentclass[11pt]{article}
% \end{verbatim}
% \end{quote}

% To load the style file in the review version:
% \begin{quote}
% \begin{verbatim}
% \usepackage[review]{acl}
% \end{verbatim}
% \end{quote}
% For the final version, omit the \verb|review| option:
% \begin{quote}
% \begin{verbatim}
% \usepackage{acl}
% \end{verbatim}
% \end{quote}

% To use Times Roman, put the following in the preamble:
% \begin{quote}
% \begin{verbatim}
% \usepackage{times}
% \end{verbatim}
% \end{quote}
% (Alternatives like txfonts or newtx are also acceptable.)

% Please see the \LaTeX{} source of this document for comments on other packages that may be useful.

% Set the title and author using \verb|\title| and \verb|\author|. Within the author list, format multiple authors using \verb|\and| and \verb|\And| and \verb|\AND|; please see the \LaTeX{} source for examples.

% By default, the box containing the title and author names is set to the minimum of 5 cm. If you need more space, include the following in the preamble:
% \begin{quote}
% \begin{verbatim}
% \setlength\titlebox{<dim>}
% \end{verbatim}
% \end{quote}
% where \verb|<dim>| is replaced with a length. Do not set this length smaller than 5 cm.


% \section{Document Body}

% \subsection{Footnotes}

% Footnotes are inserted with the \verb|\footnote| command.\footnote{This is a footnote.}

% \subsection{Tables and figures}

% See Table~\ref{tab:accents} for an example of a table and its caption.
% \textbf{Do not override the default caption sizes.}

% \begin{table}
%   \centering
%   \begin{tabular}{lc}
%     \hline
%     \textbf{Command} & \textbf{Output} \\
%     \hline
%     \verb|{\"a}|     & {\"a}           \\
%     \verb|{\^e}|     & {\^e}           \\
%     \verb|{\`i}|     & {\`i}           \\
%     \verb|{\.I}|     & {\.I}           \\
%     \verb|{\o}|      & {\o}            \\
%     \verb|{\'u}|     & {\'u}           \\
%     \verb|{\aa}|     & {\aa}           \\\hline
%   \end{tabular}
%   \begin{tabular}{lc}
%     \hline
%     \textbf{Command} & \textbf{Output} \\
%     \hline
%     \verb|{\c c}|    & {\c c}          \\
%     \verb|{\u g}|    & {\u g}          \\
%     \verb|{\l}|      & {\l}            \\
%     \verb|{\~n}|     & {\~n}           \\
%     \verb|{\H o}|    & {\H o}          \\
%     \verb|{\v r}|    & {\v r}          \\
%     \verb|{\ss}|     & {\ss}           \\
%     \hline
%   \end{tabular}
%   \caption{Example commands for accented characters, to be used in, \emph{e.g.}, Bib\TeX{} entries.}
%   \label{tab:accents}
% \end{table}

% As much as possible, fonts in figures should conform
% to the document fonts. See Figure~\ref{fig:experiments} for an example of a figure and its caption.

% Using the \verb|graphicx| package graphics files can be included within figure
% environment at an appropriate point within the text.
% The \verb|graphicx| package supports various optional arguments to control the
% appearance of the figure.
% You must include it explicitly in the \LaTeX{} preamble (after the
% \verb|\documentclass| declaration and before \verb|\begin{document}|) using
% \verb|\usepackage{graphicx}|.

% \begin{figure}[t]
%   \includegraphics[width=\columnwidth]{example-image-golden}
%   \caption{A figure with a caption that runs for more than one line.
%     Example image is usually available through the \texttt{mwe} package
%     without even mentioning it in the preamble.}
%   \label{fig:experiments}
% \end{figure}

% \begin{figure*}[t]
%   \includegraphics[width=0.48\linewidth]{example-image-a} \hfill
%   \includegraphics[width=0.48\linewidth]{example-image-b}
%   \caption {A minimal working example to demonstrate how to place
%     two images side-by-side.}
% \end{figure*}

% \subsection{Hyperlinks}

% Users of older versions of \LaTeX{} may encounter the following error during compilation:
% \begin{quote}
% \verb|\pdfendlink| ended up in different nesting level than \verb|\pdfstartlink|.
% \end{quote}
% This happens when pdf\LaTeX{} is used and a citation splits across a page boundary. The best way to fix this is to upgrade \LaTeX{} to 2018-12-01 or later.

% \subsection{Citations}

% \begin{table*}
%   \centering
%   \begin{tabular}{lll}
%     \hline
%     \textbf{Output}           & \textbf{natbib command} & \textbf{ACL only command} \\
%     \hline
%     \citep{Gusfield:97}       & \verb|\citep|           &                           \\
%     \citealp{Gusfield:97}     & \verb|\citealp|         &                           \\
%     \citet{Gusfield:97}       & \verb|\citet|           &                           \\
%     \citeyearpar{Gusfield:97} & \verb|\citeyearpar|     &                           \\
%     \citeposs{Gusfield:97}    &                         & \verb|\citeposs|          \\
%     \hline
%   \end{tabular}
%   \caption{\label{citation-guide}
%     Citation commands supported by the style file.
%     The style is based on the natbib package and supports all natbib citation commands.
%     It also supports commands defined in previous ACL style files for compatibility.
%   }
% \end{table*}

% Table~\ref{citation-guide} shows the syntax supported by the style files.
% We encourage you to use the natbib styles.
% You can use the command \verb|\citet| (cite in text) to get ``author (year)'' citations, like this citation to a paper by \citet{Gusfield:97}.
% You can use the command \verb|\citep| (cite in parentheses) to get ``(author, year)'' citations \citep{Gusfield:97}.
% You can use the command \verb|\citealp| (alternative cite without parentheses) to get ``author, year'' citations, which is useful for using citations within parentheses (e.g. \citealp{Gusfield:97}).

% A possessive citation can be made with the command \verb|\citeposs|.
% This is not a standard natbib command, so it is generally not compatible
% with other style files.

% \subsection{References}

% \nocite{Ando2005,andrew2007scalable,rasooli-tetrault-2015}

% The \LaTeX{} and Bib\TeX{} style files provided roughly follow the American Psychological Association format.
% If your own bib file is named \texttt{custom.bib}, then placing the following before any appendices in your \LaTeX{} file will generate the references section for you:
% \begin{quote}
% \begin{verbatim}
% \bibliography{custom}
% \end{verbatim}
% \end{quote}

% You can obtain the complete ACL Anthology as a Bib\TeX{} file from \url{https://aclweb.org/anthology/anthology.bib.gz}.
% To include both the Anthology and your own .bib file, use the following instead of the above.
% \begin{quote}
% \begin{verbatim}
% \bibliography{anthology,custom}
% \end{verbatim}
% \end{quote}

% Please see Section~\ref{sec:bibtex} for information on preparing Bib\TeX{} files.

% \subsection{Equations}

% An example equation is shown below:
% \begin{equation}
%   \label{eq:example}
%   A = \pi r^2
% \end{equation}

% Labels for equation numbers, sections, subsections, figures and tables
% are all defined with the \verb|\label{label}| command and cross references
% to them are made with the \verb|\ref{label}| command.

% This an example cross-reference to Equation~\ref{eq:example}.

% \subsection{Appendices}

% Use \verb|\appendix| before any appendix section to switch the section numbering over to letters. See Appendix~\ref{sec:appendix} for an example.

% \section{Bib\TeX{} Files}
% \label{sec:bibtex}

% Unicode cannot be used in Bib\TeX{} entries, and some ways of typing special characters can disrupt Bib\TeX's alphabetization. The recommended way of typing special characters is shown in Table~\ref{tab:accents}.

% Please ensure that Bib\TeX{} records contain DOIs or URLs when possible, and for all the ACL materials that you reference.
% Use the \verb|doi| field for DOIs and the \verb|url| field for URLs.
% If a Bib\TeX{} entry has a URL or DOI field, the paper title in the references section will appear as a hyperlink to the paper, using the hyperref \LaTeX{} package.

% \section*{Limitations}

% This document does not cover the content requirements for ACL or any
% other specific venue.  Check the author instructions for
% information on
% maximum page lengths, the required ``Limitations'' section,
% and so on.

% \section*{Acknowledgments}

% This document has been adapted
% by Steven Bethard, Ryan Cotterell and Rui Yan
% from the instructions for earlier ACL and NAACL proceedings, including those for
% ACL 2019 by Douwe Kiela and Ivan Vuli\'{c},
% NAACL 2019 by Stephanie Lukin and Alla Roskovskaya,
% ACL 2018 by Shay Cohen, Kevin Gimpel, and Wei Lu,
% NAACL 2018 by Margaret Mitchell and Stephanie Lukin,
% Bib\TeX{} suggestions for (NA)ACL 2017/2018 from Jason Eisner,
% ACL 2017 by Dan Gildea and Min-Yen Kan,
% NAACL 2017 by Margaret Mitchell,
% ACL 2012 by Maggie Li and Michael White,
% ACL 2010 by Jing-Shin Chang and Philipp Koehn,
% ACL 2008 by Johanna D. Moore, Simone Teufel, James Allan, and Sadaoki Furui,
% ACL 2005 by Hwee Tou Ng and Kemal Oflazer,
% ACL 2002 by Eugene Charniak and Dekang Lin,
% and earlier ACL and EACL formats written by several people, including
% John Chen, Henry S. Thompson and Donald Walker.
% Additional elements were taken from the formatting instructions of the \emph{International Joint Conference on Artificial Intelligence} and the \emph{Conference on Computer Vision and Pattern Recognition}.

% Bibliography entries for the entire Anthology, followed by custom entries
%\bibliography{anthology,custom}
% Custom bibliography entries only
\bibliography{custom}

% =====================================================================
% PREAMBLE ADDITIONS (one-time)
% =====================================================================
% No tcolorbox needed any more.
% Add to your preamble (or remove if unused elsewhere):
%   \usepackage{xcolor}     % only if not already loaded for other reasons
%
% Custom stage-block style: bold rule, bold heading, indented content.
% Survives greyscale printing because nothing depends on colour fill.

\appendix

\newcommand{\stageblock}[2]{%
  \par\vspace{0.5\baselineskip}%
  \noindent\rule{\linewidth}{0.6pt}\par
  \noindent\textbf{#1}\par
  \vspace{0.15\baselineskip}%
  \begingroup
    \footnotesize
    \setlength{\parskip}{0pt}%
    \setlength{\leftskip}{0.5em}%
    #2
  \endgroup
  \par\vspace{0.25\baselineskip}%
}

% Compact sub-group label inside a stage block.
\newcommand{\substage}[1]{\par\vspace{0.15\baselineskip}\noindent\emph{#1}\par}

% =====================================================================
\section{End-to-End Walkthrough}
\label{app:walkthrough}

To make the pipeline concrete, we trace HiReKG end-to-end on a single source paragraph. The blocks below show, in order, the source paragraph, subject extraction (Section~\ref{sec:method-subject}), the raw \texttt{(subject, relation, object)} triples produced by property extraction (Section~\ref{sec:method-property}), the detection step that classifies each triple by how its object relates to the extracted entity set, the triples after object resolution (Section~\ref{sec:method-resolution}), and the final canonical triples after refinement. We omit some triples (marked \texttt{\dots}) for space; the full trace is available with the code release.

\stageblock{Source paragraph (2 chunks)}{%
A butterfly's life cycle has four stages: egg, larva, pupa, and adult. Female butterflies lay eggs on host plants, usually on the underside of leaves. After hatching, the larva, called a caterpillar, feeds on leaves and grows rapidly, shedding its skin through a process known as molting. The caterpillar then attaches to a stem and forms a chrysalis, inside which it undergoes metamorphosis into a butterfly.\par
The adult butterfly feeds on nectar through a tubular mouthpart called a proboscis, and pollinates flowers by transferring pollen between them. Adult butterflies typically live a few days to several weeks. Females then lay eggs on host plants, beginning the cycle again.%
}

\stageblock{Stage 1 --- Subject Extraction (25 subjects)}{%
\texttt{butterfly}, \texttt{life cycle}, \texttt{egg}, \texttt{larva}, \texttt{pupa}, \texttt{adult}, \texttt{female butterflies}, \texttt{host plants}, \texttt{leaves}, \texttt{caterpillar}, \texttt{skin}, \texttt{molting}, \texttt{stem}, \texttt{chrysalis}, \texttt{metamorphosis}, \texttt{adult butterfly}, \texttt{nectar}, \texttt{tubular mouthpart}, \texttt{proboscis}, \texttt{flowers}, \texttt{pollen}, \texttt{females}, \texttt{cycle}, \dots%
}

\stageblock{Stage 2 --- Property Extraction (raw triples, 20 total)}{%
\texttt{(life cycle, has, four stages: egg, larva, pupa, and adult)}\par
\texttt{(female butterflies, lay, eggs on host plants)}\par
\texttt{(female butterflies, lay eggs on, the underside of leaves)}\par
\texttt{(larva, is called, a caterpillar)}\par
\texttt{(larva, grows, rapidly)}\par
\texttt{(larva, sheds, its skin)}\par
\texttt{(caterpillar, undergoes, metamorphosis into a butterfly)}\par
\texttt{(adult butterfly, feeds on, nectar)}\par
\texttt{(adult butterfly, feeds through, a tubular mouthpart called a proboscis)}\par
\texttt{(adult butterfly, live, a few days to several weeks)}\par
\dots%
}

\stageblock{Stage 3 --- Detection (route each triple by entity count in object)}

The object string is normalized and matched against the extracted entity set. In the displayed triples, \texttt{[prefix]} marks the text before a matched entity, \texttt{\{suffix\}} marks the text after, and \texttt{\guillemotleft entity\guillemotright} marks a matched entity span.{%


\textbf{\substage{Clean (no embedded entity, forwarded as-is)}}
\texttt{(larva, grows, rapidly)}\par
\texttt{(adult butterfly, live, a few days to several weeks)}

textbf{\substage{Single embedded entity (deterministic prefix/suffix split)}}
\texttt{(larva, is called, [a] \guillemotleft caterpillar\guillemotright)}\par
\texttt{(larva, sheds, [its] \guillemotleft skin\guillemotright)}\par
\texttt{(caterpillar, forms, [a] \guillemotleft chrysalis\guillemotright)}\par
\texttt{(female butterflies, lay eggs on, [the underside of] \guillemotleft leaves\guillemotright)}\par
\texttt{(adult butterfly, transfers, \guillemotleft pollen\guillemotright\ \{between them\})}

\textbf{\substage{Multiple embedded entities (LLM-chosen strategy)}}
\texttt{(life cycle, has, four stages: \guillemotleft egg\guillemotright, \guillemotleft larva\guillemotright, \guillemotleft pupa\guillemotright, and \guillemotleft adult\guillemotright)}\par
\texttt{(female butterflies, lay, \guillemotleft eggs\guillemotright\ on \guillemotleft host plants\guillemotright)}\par
\texttt{(adult butterfly, feeds through, a \guillemotleft tubular mouthpart\guillemotright\ called a \guillemotleft proboscis\guillemotright)}\par
\texttt{(caterpillar, undergoes, \guillemotleft metamorphosis\guillemotright\ into a \guillemotleft butterfly\guillemotright)}\par
\dots%
}

\stageblock{Stage 4 --- Object Resolution (28 triples, grouped by strategy)}{%
\textbf{\substage{Clean (forwarded unchanged)}}
\texttt{(larva, grows, rapidly)}\par
\texttt{(adult butterfly, live, a few days to several weeks)}

\textbf{\substage{Single-entity --- prefix folded into relation}}
\texttt{(female butterflies, lay eggs on the underside of, leaves)}\par
\texttt{(larva, is called, caterpillar)}\par
\texttt{(larva, sheds its, skin)}\par
\texttt{(caterpillar, attaches to, stem)}\par
\texttt{(caterpillar, forms, chrysalis)}

\textbf{\substage{Multi-entity --- \textsc{Hub} (parallel list of objects)}}
\texttt{(life cycle, has stage, egg)}\par
\texttt{(life cycle, has stage, larva)}\par
\texttt{(life cycle, has stage, pupa)}\par
\texttt{(life cycle, has stage, adult)}

\textbf{\substage{Multi-entity --- \textsc{Chain} (sequential connectives)}}
\texttt{(female butterflies, lay, eggs)}\par
\texttt{(eggs, laid on, host plants)}\par
\texttt{(caterpillar, undergoes, metamorphosis)}\par
\texttt{(metamorphosis, into, butterfly)}

\textbf{\substage{Multi-entity --- \textsc{Nested} (appositive / definitional)}}
\texttt{(adult butterfly, feeds through a, tubular mouthpart)}\par
\texttt{(tubular mouthpart, called a, proboscis)}\par
\dots%
}

\stageblock{Stage 5 --- Refinement (canonicalization, 24 entities)}

Refinement canonicalizes surface variants in the entity set. On this trace, the plural form \texttt{eggs} is folded into the singular \texttt{egg}, and the redundant second occurrence of \texttt{host plants} is removed, reducing the entity set from 25 to 24. Triples referencing the dropped surface forms are rewritten against the canonical entity, yielding the final knowledge graph.%



A brief commentary on what each stage contributes. Subject extraction recovers the entities that anchor the document. Property extraction lifts the entity-pair constraint, allowing free-form objects: this is what preserves single-entity facts such as \texttt{(larva, grows, rapidly)} and \texttt{(adult butterfly, live, a few days to several weeks)}, which entity-first pipelines drop by construction. The detection step then partitions raw triples into a clean set, a single-embedded-entity set that the prefix/suffix rule resolves deterministically, and a multi-embedded-entity set routed to the LLM for \textsc{Hub}, \textsc{Chain}, or \textsc{Nested} decomposition. Object resolution applies these branches: a single raw triple about the four life-cycle stages becomes four \texttt{has stage} edges, and the appositive \texttt{a tubular mouthpart called a proboscis} becomes a two-hop chain. Refinement performs surface-form deduplication without changing semantics.


\section{Object Resolution: Worked Examples}
\label{app:obj-res-examples}

Table~\ref{tab:obj-res-examples} shows one worked example for each of the five resolution patterns introduced in Section~\ref{sec:method-resolution}. The first four rows are drawn from the end-to-end walkthrough in Appendix~\ref{app:walkthrough}; the \textsc{Nested} row uses a constructed example, since the butterfly paragraph does not contain an appositive in the object position.

\begin{table*}[t!]
\centering
\small
\setlength{\tabcolsep}{5pt}
\renewcommand{\arraystretch}{1.2}
\begin{tabular}{p{0.13\linewidth} | p{0.27\linewidth} | p{0.13\linewidth} | p{0.34\linewidth}}
\toprule
\textbf{Pattern} & \textbf{Raw triple (Stage 2 output)} & \textbf{Embedded entities} & \textbf{Resolved triples (Stage 4 output)} \\
\midrule
Clean & \texttt{(larva, grows, rapidly)} & \emph{none} & \texttt{(larva, grows, rapidly)} \\
\midrule
Single-entity \newline (prefix) & \texttt{(female butterflies, lay eggs on, the underside of leaves)} & \texttt{leaves} & \texttt{(female butterflies, lay eggs on the underside of, leaves)} \\
\midrule
Single-entity \newline (suffix) & \texttt{(adult butterfly, transfers, pollen between them)} & \texttt{pollen} & \texttt{(adult butterfly, among themselves transfers, pollen)} \\
\midrule
\textsc{Hub} & \texttt{(life cycle, has, four stages: egg, larva, pupa, and adult)} & \texttt{egg}, \texttt{larva}, \texttt{pupa}, \texttt{adult} & \texttt{(life cycle, has stage, egg)} \newline \texttt{(life cycle, has stage, larva)} \newline \texttt{(life cycle, has stage, pupa)} \newline \texttt{(life cycle, has stage, adult)} \\
\midrule
\textsc{Chain} & \texttt{(female butterflies, lay, eggs on host plants)} & \texttt{eggs}, \texttt{host plants} & \texttt{(female butterflies, lay, eggs)} \newline \texttt{(eggs, laid on, host plants)} \\
\midrule
\textsc{Nested}\textsuperscript{$\ddagger$} & \texttt{(treaty, signed in, Paris, the capital of France)} & \texttt{Paris}, \texttt{France} & \texttt{(treaty, signed in, Paris)} \newline \texttt{(Paris, capital of, France)} \\
\bottomrule
\end{tabular}
\caption{Worked examples for the five object-resolution patterns. The first four rows come from the butterfly walkthrough (Appendix~\ref{app:walkthrough}); the \textsc{Nested} row is constructed. \textsuperscript{$\ddagger$}~Constructed example; the butterfly paragraph does not contain an appositive in the object position.}
\label{tab:obj-res-examples}
\end{table*}


\section{Evaluation Prompts and Judge Configuration}
\label{app:metrics}

Both MINE Score and Faithfulness Score are LLM-as-judge metrics. This appendix provides the verbatim prompts and the judge configuration used in all reported experiments.

\paragraph{Judge models.} To avoid same-model bias between extractor and judge, the judge is matched by tier: \textbf{GPT-5}~\citep{gpt5} judges all knowledge graphs extracted by GPT-4o, and \textbf{Qwen3-32B}~\citep{qwen3_32b_fp8} judges all knowledge graphs extracted by Qwen3-14B and Qwen3-8B. The judge model is held fixed across HiReKG and both baselines within each tier, so the comparison is symmetric. Both judges are queried at high reasoning effort, with structured JSON output, and outputs that fail to parse as valid JSON are counted as judge errors and excluded from the denominator (see scoring formulas below).

\paragraph{MINE Score (recall).} For each ground-truth atomic fact, top-$k{=}8$ KG nodes are retrieved by cosine similarity over \texttt{all-MiniLM-L6-v2} embeddings of the fact against each node name. Each retrieved node is expanded by a two-hop traversal over outgoing and incoming edges, and every traversed edge $(s, r, o)$ is serialized as the sentence \texttt{"$s$ $r$ $o$."} (or \texttt{"$o$ $r$ $s$."} for incoming edges). The resulting sentences are joined with single spaces into a flat context blob. We use the canonical retrieval routine from \texttt{kg-gen}~\citep{kggen} unchanged across all methods. The judge is then prompted with the context blob and the atomic fact:

\stageblock{MINE Score judge prompt (system + user)}{%
\textit{System:} You are an evaluator. Determine whether the provided context contains the information stated in the given correct answer. Think step-by-step internally, then respond with valid JSON of the form \texttt{\{"reasoning": "<brief reasoning>", "evaluation": 0 | 1\}}. Rules: \texttt{evaluation = 1} if the context contains (supports) the correct answer, \texttt{0} otherwise. Evaluate on information content, not exact wording; paraphrases are fine. If the context is empty or irrelevant, return \texttt{0}. Output ONLY the JSON object, no code fences, no extra keys.\par
\textit{User:} \texttt{Context:\textbackslash n---\textbackslash n\{context\}\textbackslash n---\textbackslash n\textbackslash nCorrect answer:\textbackslash n---\textbackslash n\{correct\_answer\}\textbackslash n---\textbackslash n\textbackslash nDoes the context contain the information stated in the correct answer? Respond with JSON.}%
}

\noindent The per-document MINE Score is the fraction of facts judged \texttt{evaluation = 1}, averaged across documents.

\paragraph{Faithfulness Score (precision).} For each predicted triple $(s, r, o)$, a verbalization LLM converts the triple into a single self-contained English sentence (the \emph{statement}) using the prompt below. The same retrieval embedder \texttt{all-MiniLM-L6-v2} is then used to retrieve the top-$k{=}5$ most similar chunks of the source document as context. The judge receives the context and the statement and emits a \texttt{SUPPORTED}/\texttt{NOT\_SUPPORTED} verdict with a required verbatim supporting span:

\stageblock{Faithfulness Score verbalization prompt}{%
You convert a structured knowledge-graph triple (subject, predicate, object) into ONE self-contained English sentence that asserts the same fact.\par
\smallskip
\textit{Strict rules:} Produce exactly ONE declarative sentence with no preamble. Preserve the subject and object strings verbatim; do not paraphrase, expand, abbreviate, change capitalization, or substitute synonyms. The predicate may be lightly inflected (tense, articles, auxiliary verbs) only as needed to form a grammatical sentence. Do not introduce information not present in the triple. Do not use pronouns; refer to the subject and object by their exact strings. Output only the sentence.\par
\smallskip
\textit{Input:} \texttt{subject: \{subject\}; predicate: \{predicate\}; object: \{object\}}.%
}

\stageblock{Faithfulness Score judge prompt}{%
You are an entailment judge for knowledge-graph faithfulness evaluation. Given a SOURCE TEXT and a STATEMENT (verbalized from a predicted KG triple), decide whether the source text SUPPORTS the statement.\par
\smallskip
\textit{Verdicts:} \texttt{SUPPORTED} if the source text explicitly entails the statement (paraphrases count, but the same proposition must be expressible from the source's wording). \texttt{NOT\_SUPPORTED} if the source either contradicts the statement or is simply silent on it.\par
\smallskip
\textit{Strict rules:} Use ONLY the source text; ignore world knowledge and prior beliefs. Before deciding, attempt to identify a verbatim span from the source. If a span supports or contradicts the statement, copy it exactly; if the source is silent, set \texttt{supporting\_span} to \texttt{null}. The \texttt{reasoning} field (1--2 sentences) must explicitly relate the cited span (or its absence) to the statement. Output a single strict JSON object with keys \texttt{supporting\_span}, \texttt{reasoning}, \texttt{verdict} in that order. No code fences, no extra keys.%
}

\noindent The verbalizer runs on GPT-5-mini (closed tier) or Qwen3-32B (open tier) at low reasoning effort. The per-document Faithfulness Score is
\[
\text{Faithfulness} = \frac{n_{\text{supported}}}{n_{\text{total}} - n_{\text{malformed}} - n_{\text{judge\_errors}}},
\]
where $n_{\text{malformed}}$ counts triples whose verbalization failed schema validation and $n_{\text{judge\_errors}}$ counts judge outputs that failed JSON parsing. The dataset score is the mean of per-document scores.

% \section{Implementation Details}
% \label{app:implementation}

% This appendix records the implementation details required to reproduce the results in the main paper. Section~\ref{app:models} lists the language models and inference settings, Section~\ref{app:chunking} describes the chunking and preprocessing used by each method, Section~\ref{app:thresholds} reports the thresholds used in \textbf{HiReKG} and in the evaluation pipelines, Section~\ref{app:baselines} describes the baseline implementations, Section~\ref{app:hardware} reports the serving stack and hardware, and Section~\ref{app:prompts} lists every prompt used in the system.

\section{Models and Hyperparameters}
\label{app:hardware}

All extraction and judging calls use the same model snapshots across HiReKG, KGGen, and RAKG, so observed differences reflect pipeline design rather than the choice of LLM. The extractors are \texttt{gpt-4o-2024-08-06} (matching the snapshot in \citealp{kggen}) and the FP8-quantized open-source models \texttt{Qwen/Qwen3-14B-FP8} and \texttt{Qwen/Qwen3-8B-FP8}. The judge is \texttt{gpt-5-2025-08-07} for the GPT-4o runs and \texttt{Qwen/Qwen3-32B-FP8} for the Qwen3 runs; the latter is held distinct from both Qwen3 extractors to avoid same-model bias.

\paragraph{Decoding and chunking.} Table~\ref{tab:hparams} summarizes per-method decoding and chunking. Top-$p$ and seed are left unset; streaming is disabled so responses return as complete JSON. Chunking is part of each published pipeline and is not normalized across methods.

\begin{table}[t!]
\centering
\small
\setlength{\tabcolsep}{2.5pt}
\begin{tabular}{lccc}
\toprule
\textbf{Setting} & \textbf{HiReKG} & \textbf{KGGen} & \textbf{RAKG} \\
\midrule
Temperature       & 0.0 / 0.2 & 0.0 & 0.0 \\
Max output tokens & 13{,}000  & 10{,}000 & 13{,}000 \\
Max retries       & 3         & 3        & 3 \\
\midrule
Chunk unit        & sentences & characters & sentences \\
Chunk size        & 3--4 sent.\textsuperscript{$\dagger$} & 3000 chars & 1 sent. \\
Sentencizer       & \texttt{en\_core\_web\_sm} & --- & native \\
\bottomrule
\end{tabular}
\caption{Decoding and chunking configuration. \textsuperscript{$\dagger$}~HiReKG uses 3 sentences per chunk in OpenAI Batch mode (temperature 0.0) and 4 sentences per chunk in non-batched mode (temperature 0.2).}
\label{tab:hparams}
\end{table}

\paragraph{Pipeline thresholds.} Every numerical threshold used inside HiReKG and the evaluation pipelines is listed below.

\smallskip
\noindent\textit{HiReKG internals.} Entity refinement mode \texttt{llm} (batch) with \texttt{WordNetLemmatizer} normalization. Semantic linking uses \texttt{BAAI/bge-m3} embeddings with \texttt{cluster\_distance\_threshold}$=0.99$ and \texttt{min\_cluster\_size}$=2$.

\smallskip
\noindent\textit{Evaluation pipelines.} \textbf{MINE Score} retrieves the top-$k{=}8$ KG nodes by \texttt{all-MiniLM-L6-v2} cosine similarity and expands to a 2-hop subgraph (defaults from \citealp{kggen}). \textbf{Faithfulness Score} retrieves the top-$k{=}5$ most similar source chunks for each verbalized triple, also using \texttt{all-MiniLM-L6-v2}. \textbf{UIR} embeds verbalized triples with \texttt{all-MiniLM-L6-v2} and clusters by pairwise BERTScore.

\smallskip
\noindent\textit{Baselines.} KGGen deduplication uses \texttt{semhash\_threshold}$=0.95$. RAKG retrieval uses \texttt{BAAI/bge-m3} embeddings with \texttt{similarity\_threshold}$=0.60$ and \texttt{retrieval\_top\_k}$=5$.

\subsection{Baseline Implementations}
\label{app:baselines}

We use the official open-source implementations of both baselines without modification to their core logic. 

\subsection{Serving Stack and Hardware}
\label{app:serving_stack}

The two Qwen3 models are served locally with \texttt{vLLM} using FP8 quantization. The GPT-4o and GPT-5 calls are issued through the OpenAI API; GPT-4o extraction calls additionally use the OpenAI Batch API for cost efficiency on the larger Qwen-comparable runs. Experiments are distributed across two servers: one with a single NVIDIA A100 80GB GPU and one with a single NVIDIA H100 80GB GPU. Different datasets are run on whichever server is free at the time, so the assignment of dataset to GPU type is not informative; the extraction LLMs themselves are deterministic enough at temperature $0.0$--$0.2$ that the GPU model does not materially affect the reported scores.


\section{Unique Information Ratio}
\label{app:uir}

\paragraph{Metric.} The Unique Information Ratio (UIR) measures redundancy among predicted triples. Each triple $(s, r, o)$ is verbalized into a short sentence and embedded with \texttt{all-MiniLM-L6-v2}; pairs whose embedding similarity exceeds a threshold are clustered together, so that paraphrased triples expressing the same fact collapse into a single cluster. UIR is defined as
\[
\mathrm{UIR} = \frac{n_{\text{clusters}}}{n_{\text{triples}}},
\]
where higher values indicate that every triple contributes a distinct fact and lower values indicate redundancy. A graph in which every triple is unique has $\mathrm{UIR}=1.0$; a graph whose triples collapse into half as many distinct facts has $\mathrm{UIR}=0.5$.

\paragraph{Per-dataset results.} Table~\ref{tab:uir} reports UIR for each method on each dataset, averaged across the three extractor LLMs (GPT-4o, Qwen3-14B, Qwen3-8B).


\begin{table*}[t!]
\centering
\small
\setlength{\tabcolsep}{4pt}
\renewcommand{\arraystretch}{1.15}
\begin{tabular}{l|rrr|rrr|rrr}
\toprule
& \multicolumn{3}{c|}{\textbf{MINE-1}} & \multicolumn{3}{c|}{\textbf{SciERC}} & \multicolumn{3}{c}{\textbf{Re-DocRED}} \\
\textbf{LLM} & Rep.\ \% & Fail \% & Yield \% & Rep.\ \% & Fail \% & Yield \% & Rep.\ \% & Fail \% & Yield \% \\
\midrule
GPT-4o    & 72.0 &  1.5 & $+$37.8 & 57.3 &  0.3 & $+$17.5 & 73.8 &  0.4 & $+$38.5 \\
Qwen3-14B & 63.9 & 13.4 & $+$13.3 & 68.7 & 11.7 & $+$14.5 & 74.4 &  9.5 & $+$9.1  \\
Qwen3-8B  & 64.9 &  2.4 & $+$11.1 & 68.8 &  2.2 & $+$13.2 & 66.9 &  2.4 & $+$12.2 \\
\bottomrule
\end{tabular}
\caption{Object resolution statistics across nine LLM--dataset configurations. Aggregated across configurations (triple-weighted), the module repairs $67.0\%$ of raw triples, fails on $5.6\%$, and increases the final triple count by $+20.4\%$.}
\label{tab:obj-res-stats}
\end{table*}



\begin{table}[h]
\centering
\small
\setlength{\tabcolsep}{6pt}
\renewcommand{\arraystretch}{1.15}
\begin{tabular}{l|ccc|c}
\toprule
\textbf{Method} & \textbf{MINE-1} & \textbf{SciERC} & \textbf{Re-DocRED} & \textbf{Overall} \\
\midrule
KGGen           & \textbf{0.69} & 0.80          & 0.54          & 0.68 \\
RAKG            & 0.58          & 0.69          & 0.49          & 0.59 \\
\textbf{HiReKG} & 0.65          & \textbf{0.83} & \textbf{0.65} & \textbf{0.71} \\
\bottomrule
\end{tabular}
\caption{Unique Information Ratio (UIR), averaged across GPT-4o, Qwen3-14B, and Qwen3-8B for each (method, dataset) pair. \textbf{Bold} marks the best per column.}
\label{tab:uir}
\end{table}



\paragraph{Discussion.} HiReKG achieves the highest overall UIR ($0.71$), ahead of KGGen ($0.68$) and RAKG ($0.59$). The advantage is largest on Re-DocRED ($+10.9$ pp over KGGen), where the longer Wikipedia-style documents give entity-first methods more opportunities to emit paraphrased triples, and on SciERC ($+2.9$ pp over KGGen). KGGen edges HiReKG on MINE-1 ($0.69$ vs.\ $0.65$, a $3.7$ pp gap), reflecting KGGen's strict entity-pair anchoring on a dataset where the short LLM-generated documents naturally yield fewer redundant triples to begin with. RAKG trails on every dataset, consistent with its much higher absolute triple counts (38{,}249 raw triples across all configurations, against $25{,}925$ for HiReKG and $16{,}379$ for KGGen) inflating the denominator faster than the cluster count.


\section{Inference Cost}
\label{app:cost}

Table~\ref{tab:cost} reports the average per-document token cost of each method, averaged across the three extractor LLMs (GPT-4o, Qwen3-14B, Qwen3-8B). Counts are total tokens billed to the inference API, decomposed into input and output.

\begin{table}[h]
\centering
\small
\setlength{\tabcolsep}{4pt}
\renewcommand{\arraystretch}{1.1}
\begin{tabular}{ll|rrr}
\toprule
\textbf{Method} & \textbf{Tokens} & \textbf{MINE-1} & \textbf{SciERC} & \textbf{Re-DocRED} \\
\midrule
\multirow{3}{*}{KGGen}
  & Input  &   4{,}829 &   2{,}107 &   2{,}161 \\
  & Output &  16{,}721 &   3{,}864 &   6{,}959 \\
  & Total  &  21{,}550 &   5{,}971 &   9{,}120 \\
\midrule
\multirow{3}{*}{RAKG}
  & Input  &  65{,}889 &   6{,}803 &  29{,}140 \\
  & Output &  99{,}079 &  17{,}877 &  51{,}336 \\
  & Total  & 164{,}968 &  24{,}680 &  80{,}476 \\
\midrule
\multirow{3}{*}{\textbf{HiReKG}}
  & Input  &  22{,}150 &   4{,}653 &   8{,}808 \\
  & Output &  42{,}382 &  10{,}936 &  18{,}534 \\
  & Total  &  64{,}532 &  15{,}589 &  27{,}342 \\
\bottomrule
\end{tabular}
\caption{Average tokens per document, averaged across the three extractor LLMs. Input/output split exposes which methods are retrieval-heavy (input) versus generation-heavy (output).}
\label{tab:cost}
\end{table}



HiReKG sits between the two baselines on every dataset. Averaging across datasets, HiReKG uses $2.9\times$ the tokens of KGGen ($35.8$k vs.\ $12.2$k per document) but only $0.40\times$ the tokens of RAKG ($35.8$k vs.\ $90.0$k). The per-dataset ratio against KGGen is stable ($2.6$--$3.0\times$), confirming that the cost overhead is constant across text styles rather than concentrated on any one dataset. Against RAKG the gap is widest on Re-DocRED, where HiReKG uses $0.34\times$ the tokens; this is the dataset where RAKG's per-document retrieval cost is highest. Output tokens dominate the cost of all three methods ($62$--$75\%$ of the total), so the differences in total cost track the differences in generation volume rather than retrieval volume. The output tokens also include the tokens used for thinking.

\section{Object Resolution Statistics}
\label{app:obj-res}

Table~\ref{tab:obj-res-stats} reports the operational statistics of the object resolution module across all nine LLM--dataset configurations. \textbf{Rep.\ \%} is the fraction of raw extractor triples that the module rewrites; \textbf{Fail \%} is the fraction whose rewrite fails validation and reverts to the original; \textbf{Yield \%} is the net change in triple count after resolution. The remainder of triples in each cell pass through unchanged.



The repair rate is tightly bounded across configurations ($57.3$--$74.4\%$), showing that object resolution fires at a stable rate independent of dataset or extractor scale. The clean rate (fraction of triples forwarded unchanged) ranges from $16.1\%$ to $42.4\%$, confirming that the module is adaptive: well-formed triples are forwarded as-is rather than rewritten. Yield is largest where GPT-4o produces dense compound objects that the module decomposes ($+38.5\%$ on Re-DocRED, $+37.8\%$ on MINE-1) and smallest where Qwen3-14B's higher fail rate ($9.5$--$13.4\%$) reverts more triples to their unprocessed form ($+9.1\%$ on Re-DocRED). The elevated Qwen3-14B fail rate is concentrated entirely in this configuration and does not appear  8B scale; we attribute it to occasional JSON-schema deviations in the mid-scale model's output, which the validation layer correctly catches before downstream stages.


\section{Single/Multi-Entity Classification}
\label{app:datasets}

The single-entity decomposition reported in Section~\ref{sec:single-entity} requires classifying every atomic fact in the recall ground truth as \emph{single-entity} (mentioning exactly one of the document's entities) or \emph{multi-entity} (mentioning two or more). The classifier differs by dataset because the available annotations differ.

\paragraph{SciERC and Re-DocRED.} Both datasets provide gold-annotated entity sets for each document. We apply deterministic string matching: each annotated entity is searched in the fact text with case-insensitive, whitespace-normalized, word-boundary-aware regex matching. Overlapping matches are deduplicated by retaining only the longest non-overlapping span, so that ``Microsoft Windows'' is not double-counted as ``Microsoft'' plus ``Windows''. A fact with one matched entity is labeled single-entity, with two or more matched entities multi-entity, and with zero matched entities ambiguous (the latter folded into multi-entity for analysis, matching the strictest possible single-entity definition).

\paragraph{MINE-1.} MINE-1 atomic facts do not come with annotated entity sets, so deterministic matching is not available. We classify each fact with GPT-5~\citep{gpt5} using a zero-shot prompt that asks the model to count the entities mentioned in the fact and emit one of \texttt{single-entity}, \texttt{multiple-entity}, or \texttt{ambiguous}. The prompt instructs GPT-5 to treat both proper nouns (e.g., ``Marie Curie'') and common-noun entities (e.g., ``butterfly'', ``chrysalis'') as entities, since MINE-1's LLM-generated explanatory prose contains few proper nouns. Ambiguous facts are again folded into multi-entity for analysis.

\paragraph{Validation against spaCy NER.} We validated the GPT-5 labels against spaCy NER (\texttt{en\_core\_web\_sm}), applying the same count-based classification rule to spaCy's entity spans. Across $1{,}500$ MINE-1 facts, GPT-5 and spaCy agreed on $76.8\%$ of labels (Cohen's $\kappa = 0.32$). The disagreement is asymmetric and concentrated where expected: spaCy systematically under-recognizes common-noun entities that off-the-shelf NER models are not trained to tag (``butterfly'', ``caterpillar'', ``chrysalis''), but agrees strongly with GPT-5 on facts containing proper nouns ($P{=}0.84$, $R{=}0.86$ on multi-entity facts). Our analysis uses the GPT-5 labels because they more accurately reflect MINE-1's actual entity structure; the residual disagreement on common-noun-heavy facts is the expected behavior of off-the-shelf NER rather than a labeling error.

\paragraph{Fact distribution.} Table~\ref{tab:fact-distribution} reports the resulting single-entity and multi-entity counts per dataset (ambiguous facts folded into multi-entity per the convention above). SciERC is the most balanced of the three; MINE-1 and Re-DocRED are similarly skewed toward multi-entity, with single-entity facts forming roughly a quarter of each dataset. Re-DocRED's smaller absolute count ($56$ single-entity facts) reflects its 15-document scope and is sufficient for the per-cell averages reported in Section~\ref{sec:single-entity}.


\begin{table*}[t]
\centering
\small
\setlength{\tabcolsep}{4pt}
\renewcommand{\arraystretch}{1.15}
\begin{tabular}{ll|c|c|c}
\toprule
\textbf{LLM} & \textbf{Method} & \textbf{MINE-1} & \textbf{SciERC} & \textbf{Re-DocRED} \\
\midrule
\multirow{3}{*}{GPT-4o}
  & KGGen           & $65.88 \pm 18.83$ & $66.23 \pm 18.18$ & $49.21 \pm 22.19$ \\
  & RAKG            & $74.33 \pm 16.37$ & $33.34 \pm 20.38$ & $51.43 \pm 15.40$ \\
  & \textbf{HiReKG} & $\mathbf{93.40 \pm \phantom{0}7.51}$ & $\mathbf{80.15 \pm 13.94}$ & $\mathbf{80.19 \pm 13.83}$ \\
\midrule
\multirow{3}{*}{Qwen3-14B}
  & KGGen           & $58.87 \pm 21.93$ & $53.81 \pm 20.34$ & $46.70 \pm 23.87$ \\
  & RAKG            & $61.33 \pm 17.76$ & $45.58 \pm 18.85$ & $54.79 \pm 13.22$ \\
  & \textbf{HiReKG} & $\mathbf{90.87 \pm \phantom{0}9.17}$ & $\mathbf{78.74 \pm 15.97}$ & $\mathbf{79.81 \pm 10.34}$ \\
\midrule
\multirow{3}{*}{Qwen3-8B}
  & KGGen           & $54.73 \pm 22.08$ & $52.52 \pm 18.63$ & $47.36 \pm 16.65$ \\
  & RAKG            & $73.80 \pm 15.12$ & $63.97 \pm 18.64$ & $55.14 \pm 14.61$ \\
  & \textbf{HiReKG} & $\mathbf{83.40 \pm 11.17}$ & $\mathbf{68.10 \pm 17.43}$ & $\mathbf{64.16 \pm 17.73}$ \\
\bottomrule
\end{tabular}
\caption{MINE Score (\%) reported as mean $\pm$ standard deviation across documents within each (method, dataset, LLM) configuration.}
\label{tab:variance}
\end{table*}

\begin{table}[h]
\centering
\small
\setlength{\tabcolsep}{6pt}
\renewcommand{\arraystretch}{1.15}
\begin{tabular}{l|rrr}
\toprule
\textbf{Dataset} & \textbf{Single} & \textbf{Multi} & \textbf{Total} \\
\midrule
MINE-1    & 344 (22.9\%) & 1{,}156 (77.1\%) & 1{,}500 \\
SciERC    & 577 (38.9\%) &   907  (61.1\%)  & 1{,}484 \\
Re-DocRED &  56 (25.1\%) &   167  (74.9\%)  & 223 \\
\bottomrule
\end{tabular}
\caption{Distribution of atomic facts across single-entity and multi-entity classes per dataset, after folding the small ambiguous bucket into multi-entity (Section~\ref{sec:single-entity}).}
\label{tab:fact-distribution}
\end{table}

% \FloatBarrier
% \section{Per-Document Variance}
\section{Per-Document Variance}
\label{app:variance}



Table~\ref{tab:variance} reports the per-document mean and standard deviation of MINE Score for each method, dataset, and extractor LLM. Each value is computed across the documents in the dataset's evaluation split, after running the full extraction and judging pipeline once per document.


\paragraph{Variance reduction.} Beyond the recall gains reported in the main paper, HiReKG also exhibits consistently lower per-document variance than either baseline. Averaged across the nine configurations, HiReKG's standard deviation is $13.0$, compared with $16.7$ for RAKG and $20.3$ for KGGen. The reduction is largest on MINE-1 with GPT-4o ($-8.9$ std relative to the strongest baseline) and holds in 8 of 9 configurations; the single exception is Re-DocRED with Qwen3-8B ($+3.1$ std), which is also the only configuration where HiReKG's mean score does not exceed the strongest baseline by a wide margin. Lower variance means HiReKG's recall advantage is distributed across documents rather than driven by a small number of high-scoring outliers, which makes the headline gains in Section~\ref{sec:main-results} more reliable as a per-document expectation.

\paragraph{What is and is not measured.} The standard deviations above capture variation \emph{across documents} within a single run of each pipeline. They do not capture variation \emph{across LLM rollouts} on the same document, which would require multi-seed runs that were not feasible under our compute budget. We note this as a limitation.
% \subsection{Code Release}
% \label{app:code}

% The full source code for HiReKG, the wrapper code for the baselines, and the evaluation pipeline will be released on publication under an open-source license, together with the prompts listed in Section~\ref{app:prompts} and the configuration files needed to reproduce every number in the main paper.

% \subsection{Prompts}
% \label{app:prompts}

% \textit{<<Prompts>>}


% \label{sec:appendix}

% \subsection{Dataset Details}
% \label{app:datasets}

% We evaluate on three datasets that together span generic, scientific, and technical domains, and that all expose ground-truth atomic facts or triples against which extracted graphs can be compared. 
% \paragraph{MINE-1.}
% MINE-1 is the recall benchmark released with KGGen \citep{kggen}. It consists of $100$ LLM-generated articles with a mean length of approximately $590$ words, spanning Arts and Society, Science, Technology, Psychology, and History. Each article is paired with $15$ manually verified atomic facts that are known to be present in the text, yielding $100 \times 15 = 1{,}500$ ground-truth facts in total. Facts are scoped to a single article and are not de-duplicated across articles, since each fact is anchored to its own source text; the recall measurement is therefore per-article and independent of overlap across articles. We use the dataset as released, without any preprocessing.

% \paragraph{SciERC (100-document subset).}
% SciERC \citep{scierc} is a relation extraction dataset of $500$ abstracts from artificial intelligence conference and workshop proceedings, annotated with six entity types and seven relation types together with coreference clusters. We use a randomly sampled subset of $100$ abstracts, drawn from the full $500$-document collection. The sampled subset contains $983$ ground-truth triples in aggregate, against which the predicted graphs are compared structurally. The abstract-level granularity provides an expert-annotated technical-domain reference that contrasts with the generic text in MINE-1.

% \paragraph{Re-DocRED (Windows-relevant subset).}
% Re-DocRED \citep{redocred} is a re-annotated version of the DocRED document-level relation extraction dataset \citep{docred} that supplements the original annotations with previously missed triples. Following \citet{lkd-kgc}, we use the subset of Re-DocRED documents that were re-annotated from scratch and that pertain to the Windows operating system domain, which provides a longer document-level setting and a denser ground-truth graph per document than the abstract-level SciERC. The subset contains 15 documents and $517$ ground-truth triples in aggregate.

% \subsection{Evaluation Metrics}
% \label{app:metrics}

% This subsection provides the full operational definitions of the four evaluation metrics used in the main paper, together with the formulas and the edge-case handling that the implementation applies. The first two metrics (MINE Score and Faithfulness Score) are treated as primary KPIs because they directly measure how much of the source-text information is captured (recall) and how much of the predicted graph is grounded in the source (precision). The remaining two (G-BERTScore and UIR) are reported as supporting structural evidence.

% \paragraph{MINE Score (recall).}
% The MINE Score measures the fraction of ground-truth atomic facts in MINE-1 that can be recovered from a predicted knowledge graph. For each article, the dataset provides $15$ manually verified facts. To score a predicted graph $\mathrm{KG}$ against an article, each fact $f$ and each node $v \in V(\mathrm{KG})$ are embedded with \texttt{all-MiniLM-L6-v2}; the top $k=8$ nodes most similar to $f$ are retrieved, and the retrieved set is expanded to include all nodes within two hops in $\mathrm{KG}$. The induced subgraph is passed to an LLM judge, which decides whether $f$ can be inferred from the subgraph. The per-article MINE Score is the percentage of facts that the judge marks as recoverable, and the dataset-level score is the mean across the $100$ articles. We extend this metric beyond MINE-1 by treating the SciERC and Re-DocRED ground-truth triples as atomic facts and applying the same retrieval-and-judgment pipeline; the resulting scores are reported under the same label in the main paper, which we make explicit here to avoid ambiguity.

% \paragraph{Faithfulness Score (precision).}
% The Faithfulness Score measures the fraction of predicted triples that are actually supported by the source text and serves as the precision counterpart to the MINE Score. The pipeline operates on each (source text, predicted triples) pair in three steps. First, each predicted triple $(s, p, o)$ is verbalized into a self-contained English statement using an LLM, with the exact entity names preserved. Second, the full source text is supplied to the judge alongside the verbalized statement; no retrieval step is used, so the judge sees the entire document. Third, the judge is required to quote the supporting span before issuing a binary verdict in $\{\textsc{Supported}, \textsc{NotSupported}\}$, which forces grounded decisions and improves auditability. The per-document Faithfulness Score is
% \begin{equation*}
% \mathrm{Faithfulness} \;=\; \frac{n_{\mathrm{supported}}}{n_{\mathrm{total}} - n_{\mathrm{malformed}} - n_{\mathrm{judge\_errors}}},
% \end{equation*}
% where $n_{\mathrm{malformed}}$ counts triples that fail verbalization (e.g.\ missing fields) and $n_{\mathrm{judge\_errors}}$ counts judgments whose output cannot be parsed; when the denominator is zero, the document is excluded from the average. The dataset-level score is reported as the mean over per-document scores, rather than a micro-average that would bias the score toward larger graphs.

% \paragraph{G-BERTScore (structural similarity).}
% G-BERTScore \citep{saha2021gbs} is a graph-level extension of BERTScore that treats each triple in a knowledge graph as a short sentence. The predicted graph $\mathrm{KG}$ and the reference graph $\mathrm{KG}^{*}$ are matched at the triple level by computing BERTScore between every pair of predicted and reference triples, after which a maximum-weight bipartite matching $M \subseteq V(\mathrm{KG}) \times V(\mathrm{KG}^{*})$ is solved to assign at most one predicted triple to each reference triple and vice versa. From $M$, graph-level precision, recall, and F1 are derived as
% \begin{equation*}
% \mathrm{G\text{-}BS\text{-}P} = \frac{\sum_{(t,t^{*}) \in M} \mathrm{BS}(t, t^{*})}{|\mathrm{KG}|}, \quad
% \mathrm{G\text{-}BS\text{-}R} = \frac{\sum_{(t,t^{*}) \in M} \mathrm{BS}(t, t^{*})}{|\mathrm{KG}^{*}|},
% \end{equation*}
% with G-BS-F1 their harmonic mean. We use \texttt{roberta-large} as the BERTScore backbone, which is the default of the reference \texttt{bert-score} implementation. G-BERTScore captures structural overlap between the predicted and reference graphs, but it is sensitive to surface-form differences and does not directly measure whether triples are grounded in the source text; for this reason we report it as supporting evidence rather than as a standalone judgment of method quality.

% \subsection{Unique Information Ratio (UIR)}
% \label{app:uir}

% The \textbf{Unique Information Ratio (UIR)} measures the fraction of triples in an extracted graph that carry non-overlapping information. We cluster semantically equivalent triples using BERTScore similarity above a fixed threshold, and define
% \[
% \mathrm{UIR} = \frac{\text{number of clusters}}{\text{number of triples}}.
% \]
% A UIR of $1.0$ indicates that every triple expresses a distinct fact, while lower values indicate higher redundancy: for example, a graph in which four triples collapse into two clusters has a UIR of $0.5$. The metric is complementary to Faithfulness; a graph can be highly faithful yet redundant if the same fact is restated across many surface forms.

% Table~\ref{tab:uir} reports UIR alongside the total number of predicted triples and the number of LLM calls per document on SciERC and Re-DocRED with GPT-4o as the extraction model. HiReKG matches KGGen on UIR ($0.628$ vs.\ $0.627$ on SciERC; $0.729$ vs.\ $0.738$ on Re-DocRED) while producing a comparable number of triples, indicating that the recall gains reported in the main paper are not driven by inflating the graph with paraphrased or duplicated facts. RAKG, by contrast, has substantially lower UIR ($0.45$--$0.49$) and produces $2$--$3\times$ more triples than KGGen or HiReKG, consistent with its known redundancy behavior.

% \begin{table}[h]
% \centering
% \small
% \setlength{\tabcolsep}{4pt}
% \begin{tabular}{l|ccc|ccc}
% \toprule
% & \multicolumn{3}{c|}{\textbf{SciERC}} & \multicolumn{3}{c}{\textbf{Re-DocRED}} \\
% \textbf{Method} & \#Triples & UIR & LLM Calls & \#Triples & UIR & LLM Calls \\
% \midrule
% KGGen           & 1762 & 0.6269 & 443  & 469  & \textbf{0.7381} & 74  \\
% RAKG            & 3839 & 0.4505 & 1290 & 1292 & 0.4921          & 660 \\
% \textbf{HiReKG} & 1858 & \textbf{0.6279} & 880 & 481 & 0.7289 & 162 \\
% \bottomrule
% \end{tabular}
% \caption{Number of predicted triples, UIR (higher = less redundant), and number of LLM calls on SciERC and Re-DocRED with GPT-4o as the extraction LLM. Ground-truth KGs contain 983 triples (SciERC) and 517 triples (Re-DocRED). \textbf{Bold} marks the best UIR per dataset.}
% \label{tab:uir}
% \end{table}

\end{document}
